% !TEX root = chapter2_reorg.tex
%% REORGANIZED COMPARISON DRAFT -- assembled body (math mined verbatim from thesis sources).

\chapter{Ramification After Blowup Is Tame}

Throughout this chapter, let $C$ be a smooth, projective, geometrically connected
curve over $k$. Fix a modulus

\[
    \mdls = \sum_i n_i P_i
\]

on $C$, and set $U = C \setminus \mdls$.

For each single-point submodulus $\ndls = n_i P_i$ of $\mdls$, set
$d_\ndls = \deg(\ndls)$. The effective divisor $\ndls$ determines a
$k$-rational closed point

\[
    Z_\ndls \hookrightarrow C^{(d_\ndls)}
        = \operatorname{Sym}^{d_\ndls}_k(C).
\]

Let

\[
    \pi_\ndls \colon
    X_\ndls := \operatorname{Bl}_{Z_\ndls}\bigl(C^{(d_\ndls)}\bigr)
    \longrightarrow C^{(d_\ndls)}
\]

be the blowup at this point. We denote its exceptional divisor by

\[
    E_\ndls := Z_\ndls \times_{C^{(d_\ndls)}} X_\ndls
\]

By
\Cref{prop:relative_symmetric_power_smooth,prop:exceptional_divisor_smooth_point},
$E_\ndls\cong\mathbb P_k^{d_\ndls-1}$; in particular, it is irreducible and
has codimension $1$ in $X_\ndls$. We denote its unique generic point by
$\eta_\ndls$, so that $\overline{\{\eta_\ndls\}}=E_\ndls$. Thus we have the
Cartesian diagram

\begin{equation}\label{equation:blowup_diagram_ndls:reorg}
\begin{tikzcd}
E_\ndls = \overline{\{\eta_\ndls\}}
    \arrow[r, hook] \arrow[d]
  & X_\ndls \arrow[d, "\pi_\ndls"] \\
Z_\ndls \arrow[r, hook]
  & C^{(d_\ndls)}.
\end{tikzcd}
\end{equation}

The purpose of this chapter is to prove
\Cref{theorem:SymmetricPowerOfSheavesIsTamelyRamifiedReduction}. Let $\cF$ be
the rank-one $\Lambda$-local system appearing in that theorem, and let

\[
    \cP = \operatorname{Fr}(\cF)
\]

be its $G = \Lambda^\times$-torsor of frames. By
\Cref{prop:torsor_module_equivalence,definition:symmetric_power_local_system},
there is a canonical identification

\[
    \operatorname{Fr}\bigl(\cF^{[d_\ndls]}\bigr)
        \cong \cP^{[d_\ndls]}.
\]

In particular, the two objects have the same ramification. It therefore
suffices to prove the following torsor formulation.

\begin{theorem}\label{theorem:torsors_after_blowup_is_tame:reorg}
Let $G$ be a finite abelian group, and let $\cP \to U$ be a $G$-torsor whose
ramification is bounded by $\mdls$. For every single-point submodulus
$\ndls = n_iP_i$ of $\mdls$, the symmetric-power torsor
$\cP^{[d_\ndls]}$ is tamely ramified at $\eta_\ndls$.
\end{theorem}

\medskip
\noindent\textbf{Organization of this chapter.}
The chapter is arranged to make one point visible: the three cases
$G = \Z/e\Z$ ($p \nmid e$), $G = \Z/p\Z$ and $G = \Z/p^r\Z$ are one and the
same computation, run through Kummer, Artin--Schreier, and
Artin--Schreier--Witt theory respectively.
\begin{itemize}
  \item The \emph{prelude} collects the background each case needs, in the
  order in which it is used: first the ramification theory of discrete
  valuation rings; then the ramification of $G$-torsors (the definitions that
  phrase everything below); and finally the three cyclic-extension theories,
  laid out in parallel behind a dictionary table
  (\Cref{table:cyclic_theory_dictionary:reorg}). None of it is specific to the
  present theorem.
  \item The remaining sections carry out the calculation. We first reduce
  \Cref{theorem:torsors_after_blowup_is_tame:reorg} to an explicit torsor on
  $\bG_m \subset \bP^1_k$; we then run the symmetric-power computation in all
  three cases at once, following the shared four-step skeleton (Steps A--D)
  recorded in \Cref{table:cyclic_computation_dictionary:reorg}; and we assemble
  the general finite abelian case by decomposing $G$ into its primary cyclic
  factors. A final section records the behaviour of ramification under products
  of blowups, used in \Cref{section:gcft}.
\end{itemize}


\section{Prelude: Local Ramification and Cyclic Extensions}

\subsection{Ramification of discrete valuation rings}
This section reviews the ramification of discrete valuations, beginning with the general theory of \textit{discrete valuation rings} and their 
integral closures in finite separable extensions. 
We then specialize to complete rings in Galois extensions to describe the ramification filtration via lower and upper numbering, 
following the treatments in \stackstag{0EXQ} and \cite{serreLF}.

\begin{definition}[\stackstag{09E4}]\label{definition:weakly_unramified_and_tamely_ramified_DVR:reorg}
    We say that $A \to B$ or $A \subset B$ is an extension of discrete valuation rings if $A$ and $B$ are discrete valuation rings and $A \to B$ is injective and local. In particular, if $\pi _ A$ and $\pi _ B$ are uniformizers of $A$ and $B$, then $\pi _ A = u \pi _ B^ e$ for some $e \geq 1$ and unit $u$ of $B$. The integer $e$ does not depend on the choice of the uniformizers as it is also the unique integer $\geq 1$ such that
\[ \mathfrak m_ A B = \mathfrak m_ B^ e \]
The integer $e$ is called the \textit{ramification index} of $B$ over $A$. We say that $B$ is \textit{weakly unramified} over $A$ if $e = 1$. 
If the extension of residue fields $\kappa _ A = A/\mathfrak m_ A \subset \kappa _ B = B/\mathfrak m_ B$ is finite, 
then we set $f = [\kappa _ B : \kappa _ A]$ and we call it the \textit{residual degree} or residue degree of the extension $A \subset B$.

Note that we do not require the extension of fraction fields to be finite.
\end{definition}

Now let $A$ be a discrete valuation ring with fraction field $K$, let  $L/K$ be a finite separable field extension and 
let $B \subset L$ be the integral closure of $A$ in $L$.
% \[ \xymatrix{ B \ar[r] & L \\ A \ar[u] \ar[r] & K \ar[u] } \]
Then the ring extension $A \subset B$ is finite, hence $B$ is Noetherian. 
The dimension of $B$ is $1$, hence $B$ is a Dedekind domain.
Let $\mathfrak m_1, \ldots , \mathfrak m_ n$ be the maximal ideals of $B$ (i.e., the primes lying over $\mathfrak m_ A$). 
We obtain extensions of discrete valuation rings

\[ A \subset B_{\mathfrak m_ i} \]
and hence ramification indices $e_ i$ and residue degrees $f_ i$. We have

\[ [L : K] = \sum \nolimits _{i = 1, \ldots , n} e_ i f_ i \]
Note that if $A$ is henselian (e.g. $A$ is complete) then $n=1$. 

\begin{definition}[\stackstag{09E9}]\label{definition:tamely_ramified_dvrs:reorg}
Let $A$ be a discrete valuation ring with fraction field $K$. Let $L/K$ be a finite separable extension. With $B$ and $\mathfrak m_ i$, $i = 1, \ldots , n$ 
as above, we say the extension $L/K$ is
\begin{enumerate}
    \item unramified with respect to $A$ if $e_ i = 1$ and the extension $\kappa (\mathfrak m_ i)/\kappa _ A$ is separable for all $i$,
    \item tamely ramified with respect to $A$ if either the characteristic of $\kappa _ A$ is $0$ or the characteristic of $\kappa _ A$ is $p > 0$, the field extensions $\kappa (\mathfrak m_ i)/\kappa _ A$ are separable, and the ramification indices $e_ i$ are prime to $p$, and
    \item totally ramified with respect to $A$ if $n = 1$ and the residue field extension $\kappa (\mathfrak m_1)/\kappa _ A$ is trivial.
\end{enumerate}
If the discrete valuation ring $A$ is clear from context, then we sometimes say $L/K$ is unramified, totally ramified, or tamely ramified for short.
\end{definition}

% - %
Now and for the rest of the section assume $A$ is a \textit{complete discrete valuation ring} with uniformizer $\pi$ 
and residue field $\kappa$, which we assume to be perfect. When $A$ and $\kappa$ are of the same characteristic $p > 0$, 
then $A$ contains a coefficient field $k \cong \kappa$ and a well known structure theorem holds: $A = k[[\pi ]] \cong k[[t]]$.
Let $K$ be the fraction field of $A$, then $K=k((\pi))$.

For the remainder of this section, assume $A$ is a \textit{complete discrete valuation ring} with uniformizer $\pi$ and a perfect residue 
field $\kappa$. In the equal characteristic case, where $\text{char}(A) = \text{char}(\kappa) = p > 0$, the 
Cohen Structure Theorem implies that $A$ contains a coefficient field $k \cong \kappa$, such that $A \cong k[[\pi]] \cong k[[t]]$. 
Consequently, the fraction field is given by $K = k((\pi))$.
Let $L/K$ be a finite Galois extension with Galois group $G$. 
The integral closure $B$ of $A$ in $L$ is itself a complete discrete valuation ring. 
Since the residue field $\kappa$ is perfect, the extension of residue fields is separable, and there exists a uniformizer 
$\Pi \in B$ such that $B = A[\Pi]$.

The \textit{lower numbering ramification filtration} of $G$, denoted by $(G_i)_{i \geq -1}$, is defined as:
\[ 
G_i = \{ \sigma \in G \mid v_B(\sigma (x) - x) \geq i + 1 \text{ for all } x \in B \} 
\]
where $v_B$ is the valuation on $L$ associated with $B$. In this indexing, $G_{-1} = G$ and $G_0$ 
corresponds to the \textit{inertia group} of the extension $L/K$. Note that $L/K$ is unramified if and only if 
$G_0 = \{1\}$, and it is tamely ramified if and only if $G_1 = \{1\}$.

A standard result in the theory of local fields(complete local fields with a discrete valuation and perfect residue field) 
ensures that the condition in the definition of $G_i$ need only be checked for the uniformizer $\Pi$ of $B$. 
Specifically, if we define the function
\[ 
i^L_K(\sigma) = v_B(\sigma(\Pi) - \Pi) 
\]
for $\sigma \in G$, then $G_i = \{ \sigma \in G \mid i^L_K(\sigma) \geq i + 1 \}$. 
The subgroups $G_i$ are normal in $G$ and become trivial for sufficiently large $i$.

Consider a tower of fields $K \subset E \subset L$, and let $H = \text{Gal}(L/E) \subset G$. The lower numbering is compatible with subgroups:
\[ 
G_i \cap H = H_i \quad \text{for all } i \geq -1 
\]
which follows from the identity $i^L_E = i^L_K|_{H}$. 
However, the lower numbering does not behave as simply with respect to quotients. If $H$ is normal in $G$, the image of $G_i$ in $G/H$ is generally not $(G/H)_i$. 
Instead, we have $G_i H / H = (G/H)_j$, where the index $j$ is given by the formula:
\[ 
j = \frac{1}{e_{L/E}} \sum_{\tau \in H} \min(i^L_K(\tau), i+1) - 1 
\]


In the literature, one reindexes the ramification groups by defining the Herbrand function $\phi_{L/K} : [-1, \infty ) \to [-1, \infty )$:
\[ \phi_{L/K}(i) = \frac{1}{e_{L/K}}\sum_{\sigma \in G} \min(i^L_K(\sigma), i+1)-1 = \int _0^i \frac{1}{[G_0 : G_t]} dt \]

This function is continuous, strictly increasing, and piecewise linear, making it a bijection. 
It satisfies the composition law $\phi_{L/K} = \phi_{E/K} \circ \phi_{L/E}$ for a tower of extensions $K \subset E \subset L$. 
Using this function, we define the \textit{upper numbering ramification groups} as:
\[ 
G^i = G_{\phi_{L/K}^{-1}(i)} 
\]
The primary advantage of this reindexing is its compatibility with quotients: for any normal subgroup $H \subset G$, we have:
\[ 
G^i H / H = (G/H)^i 
\]
for all $i \geq -1$.


%% [REVIEW: never used -- the structure of G_i/G_{i+1} is not invoked anywhere (tameness is always derived via Kummer/Artin-Schreier); consider deleting or reducing to a one-sentence remark]
\begin{prop}[\Cite{serreLF}, IV Corollary 3]
    the quotients $G_i/G_{i+1}$, $i \geq 1$, are abelian groups, and are direct products of cyclic groups of order $p$. 
    The group $G_1$ is a $p$-group.
\end{prop}


%% [REVIEW: never used -- same as the proposition above; nothing downstream uses the semi-direct product structure of G_0]
\begin{prop}[\Cite{serreLF}, IV Corollary 4]
    The inertia group $G_0$ has the following property:
    \textit{It is the semi-direct product of a cyclic group of order prime to $p$ with a normal subgroup whose order is a power of $p$.}
\end{prop}

% We deduce the following lemma, that will be helpful later:
% \begin{lemma}
% Let $L/K$ be finite galois extesntion with galois group $G=\Z/p^r\Z$
% Assume there exists intermediate local field $K < M < L$ such that $M/K$ is unramified.
% Then $L/K$ is unramified.
% \end{lemma}






\subsection{Ramification of \texorpdfstring{$G$}{G}-torsors}

Generally speaking, assume we are given a locally Noetherian scheme $X$, 
a dense open $U \subset X$, a finite \'etale morphsm $f: Y \to U$ and a prime divisor $Z \subset X$ such that $Z \cap U = \emptyset$
and the local ring $\Oo_{X, \xi}$ of $X$ at the generic point $\xi$ of $Z$ is a discrete valuation ring. 
Setting $K_\xi = \Frac(\Oo_{X, \xi})$ we obtain a cartesian square
\[ \xymatrix{ \mathop{\mathrm{Spec}}(K_\xi ) \ar[r] \ar[d] & U \ar[d] \\ \mathop{\mathrm{Spec}}(\mathcal{O}_{X, \xi }) \ar[r] & X } \]
of schemes. 
In particular, we see that $Y \times _ U \mathop{\mathrm{Spec}}(K_\xi )$ is the spectrum of a finite separable algebra $L_\xi /K_\xi $. 
\begin{definition}\label{definition:tamely_ramified_in_codim_1:reorg}
We say: 
\begin{enumerate}
    \item $Y$ is \textit{unramified} over $X$ at $(Z, \xi)$ resp. $Y$ is \textit{tamely ramified} over $X$ at $(Z, \xi)$
        if $L_\xi /K_\xi $ is unramified, resp. tamely ramified with respect to $\mathcal{O}_{X, \xi }$
        (\Cref{definition:tamely_ramified_dvrs:reorg})
    
    \item $Y$ is unramified over $X$ in codimension $1$, resp. $Y$ is tamely ramified over $X$ in codimension $1$ if $L_\xi /K_\xi $ is unramified, resp. tamely ramified 
    with respect to $\mathcal{O}_{X, \xi }$ for every $(Z, \xi )$ as above. 
\end{enumerate}

More precisely, we decompose $L_\xi $ into a product of finite separable field extensions of $K_\xi $ and we require each of these to be unramified, resp. tamely ramified with 
respect to $\mathcal{O}_{X, \xi }$.

\end{definition}

Now back to our orginial assumptions, $C$ is projective smooth geometrically connected curve
over a field $k$ and $U=C\setminus \mdls$. Every $(Z, \xi)$ is now a closed point $P$ and $\Oo_{C, P}$ is a discrete valuation ring. 
Assume $f: Y \to U$ is actually $G$-torsor $\cP \to U$ in $U_{\text{\'Et}}$ for $G$ finite abelian group.

By passing to the completion $\widehat{\Oo}_{C,P}$, we obtain the complete local field $\widehat{K}_P$.
Restricting the torsor $\cP$ to the punctured formal neighborhood $\Spec(\widehat{K}_P)$
 allows us to study the ramification in a strictly local context. 
 This restriction yields a $G$-torsor over a field, which necessarily decomposes into a disjoint union of spectra of finite separable field extensions
 $$ \cP|_{\Spec(\widehat{K}_P)} \cong \bigsqcup \Spec(F) $$
 such that:
 \begin{enumerate}
    \item These extensions $F/\widehat{K}_P$ are all isomorphic and Galois. 
    \item The Galois group $H = \Gal(F/\widehat{K}_P)$ identifies with the stabilizer of any connected component of the pullback $\cP|_{\Spec(\widehat{K}_P)}$, 
     which is a subgroup of $G$. 
    \item  The number of such components is then given by the index $[G:H]$.
 \end{enumerate}

 \begin{definition}\label{definition:local_ramification_boundness:reorg}
    We say that the extension $F/\widehat{K}_P$ \textit{has ramification bounded by $r$} if the ramification group $H^r$ is trivial. 
    We say \textit{$\cP$ has ramification bounded by $r$ at $P$} if the corresponding local field extensions $F/\widehat{K}_P$  
    satisfy this condition.
 \end{definition}


\begin{definition}
    Let $\mdls = \sum n_i P_i$ be a modulus on $C$. We say the $G$-torsor $\cP \to U$ has 
    \textbf{ramification bounded by $\mdls$} if, for each point $P_i$, the local restriction 
    $\cP|_{\Spec(\widehat{K}_{P_i})}$ has ramification bounded by the coefficient $n_i$.
\end{definition}

%%%%%%%%%%%%%%%%%%%%%%%%%%%%%%%%%%%%%%%
This local property remains invariant under base change to the separable closure.
Let $\bar{s} = \Spec(\bar{k}) \to \Spec(k)$ 
be a geometric point. 
The higher ramification groups for $\cP|_{\Spec(\widehat{K}_P)}$ are isomorphic to those of 
the geometric restriction $\cP_{\bar{k}}$ at any point $\bar{x}$ lying over $P$. 
Thus an equivalent defintion:
\begin{definition}
    Let $\mdls = \sum n_i P_i$ be a modulus on $C$. 
    The $G$-torsor $\cP \to U$ has \textbf{ramification bounded by $\mdls$} 
    if for every geometric point $\bar{x}$ of $\mdls$ (with image $\bar{s}$ in $\Spec(k)$), 
    the restriction of $\cP$ to$$ \Spec(\widehat{\mathcal{O}}_{C_{\bar{k}},\bar{x}}) \times_{C_{\bar{k}}} U_{\bar{k}} $$
    has ramification bounded by the multiplicity of $\mdls_{\bar{s}}$ at $\bar{x}$ in the sense of 
    \Cref{definition:local_ramification_boundness:reorg}.
\end{definition}
These definitions are equivalent because $\widehat{\mathcal{O}}_{C_{\bar{k}},\bar{x}} \cong \widehat{\mathcal{O}}_{C,P} \otimes_k \bar{k}$. Furthermore, the notions of tame ramification and unramifiedness derived here coincide with the codimension-$1$ criteria in \Cref{definition:tamely_ramified_in_codim_1:reorg}.

There is another equiviliant formulation through the language of characters.
The group $G$, being a finite abelian group over $k$, is \'etale. 
Consequently, the correspondence between $G$-torsors and the fundamental group allows us to describe $G$-torsors
 via characters. Specifically, there is an isomorphism: (\Cref{cor:abelian_equivilance_fundamental_torsors})
\[ \mathrm{Hom}_{\mathrm{cont.}}(\pi_1^{\text{\'et}}(U, \overline{x}), G) \xrightarrow{\cong} \mathrm{Tors}(U_{\text{\'et}}, G) \]

In our local setting, where we consider the restriction to the complete local field $\widehat{K}_P$, 
the fundamental group identifies with the absolute Galois group 
$G_{\widehat{K}_P} := \mathrm{Gal}({\widehat{K}_P}^{\mathrm{sep}}/{\widehat{K}_P})$. 
Thus, the restricted $G$-torsor $\mathcal{P}|_{\mathrm{Spec}(\widehat{K}_P)}$ corresponds to a unique continuous homomorphism:
\[ \rho : G_{\widehat{K}_P} \to G \]

Under this correspondence, the torsor $\mathcal{P}$ has ramification bounded by $r$ at $P$ 
if and only if the homomorphism $\rho$ trivializes the $r$-th ramification group in the upper numbering:
\[ \rho(G_{\widehat{K}_P}^r) = \{ 1 \} \]

\subsection*{Basic Properties of Ramification of \texorpdfstring{$G$-Torsors}{G-Torsors}}
We now prove some elementary lemmas regarding the ramification of $G$-torsors.
\begin{lemma}\label{lemma:bounding_ramification_of_contracted_product:reorg}
    Let $G$ be a finite abelian group and $X$ be a locally Noetherian scheme over a field $k$. 
    Let $U \subset X$ be a dense open subset and let $(Z, \xi)$ be a prime divisor in the complement $X \setminus U$.
    Assume $\cP_1$ and $\cP_2$ are two $G$-torsors on $U_{\text{ét}}$, Such that $\cP_1$ has ramification bounded by $r_1$ at $(Z, \xi)$, and $\cP_2$ has ramification bounded by $r_2$ at $(Z, \xi)$.
    Then their product $\cP_1 \otimes^G \cP_2$ has ramification bounded by $max(r_1, r_2)$ at $(Z, \xi)$.
\end{lemma}
\begin{proof}
    Let $A=\widehat{\Oo_{X, \xi}}$ and let $K=Frac(A)$.
    Let $\rho_1, \rho_2: G_K \to G$ be the associated continious homomorphisms corresponding to the $G$-torsors 
    $\cP_1|_{\Spec(K)}$, $\cP_2|_{\Spec(K)}$.
    We finish by noticing that the associated character of $(\cP_1 \otimes^G \cP_2)|_{\Spec(K)}$ is $\rho=\rho_1 + \rho_2$.    
\end{proof}

\begin{lemma}\label{lemma:descend_ramification_along_etale:reorg}

Let $f: X \to Y$ be a morphism of locally Noetherian schemes. Let $U_X \subset X$ and $U_Y \subset Y$ be dense open subschemes such that $f^{-1}(U_Y) \subset U_X$. 

Let $(Z_X, \xi_X)$ and $(Z_Y, \xi_Y)$ be prime divisors of $X$ and $Y$ such that $Z_X \cap U_X = \emptyset$ and $Z_Y \cap U_Y = \emptyset$. 
Suppose $f(\xi_X) = \xi_Y$ and that $f$ is étale at $\xi_X$.
Let $\mathcal{P}$ be a $G$-torsor over $U_Y$, and let $f^{-1} \mathcal{P}$ be its pullback to $f^{-1}(U_Y) \subset U_X$. 
Then, the ramification of $f^{-1} \mathcal{P}$ at $\xi_X$ is bounded by $r$ if and only if the ramification of $\mathcal{P}$ at $\xi_Y$ is bounded by $r$.
\end{lemma}
\begin{proof}
The boundedness of ramification is determined by the behavior of the torsor over the completion of the local rings at the generic points. 
Let $A = \mathcal{O}_{Y, \xi}$ and $B = \mathcal{O}_{X, \xi_X}$ be the discrete valuation rings at the generic points, with fraction fields 
$K$ and $L$ respectively. Since $f$ is étale at $\xi_X$, the map $A \to B$ is a flat, unramified (hence weakly unramified) local homomorphism. 
Consequently, the extension of completions $\widehat{L} / \widehat{K}$ is a finite unramified extension of complete discretely valued fields.
We finish by noticing that the upper numbering filtration on the absolute Galois group is compatible with unramified base change.\footnote{
 Specifically, let $G_K = \operatorname{Gal}(K^{sep}/K)$ and $G_L = \operatorname{Gal}(L^{sep}/L)$. For an unramified extension, the Herbrand function is the identity, which implies that for any $r \geq 0$:
\[ G_L^r = G_K^r \cap G_L \]}
\end{proof}


Let $C$ be a smooth, proper (projective) curve over a field $k$, and let $U \subset C$ be a dense open subset.
Let $G$ be a finite abelian group. and let $\pi_U: U' \to U$ be a $G$-torsor in $U_{\text{\'Et}}$. 
Let $C'$ be the normal proper model of $U'$ over $k$ (i.e. the unique smooth proper curve which has $U'$ as a dense open subset).
Let $\pi: C' \to C$ be the induced morphism. (which is finite generically \'etale over $U$). 
Let $U'^{(d)}$ be the standard $d$-th symmetric product of $U'$ as a scheme, and let $U'^{[d]}$ be the torsor-theoetric symmetric product,
by \Cref{cor:sym-power-quotient} we have a canonical morphism $q_U: U'^{(d)} \to U'^{[d]}$.
Note that $C'^{(d)}$ is smooth and finite over $C^{(d)}$, hence it is the normalization of $C^{(d)}$ inside $K(C'^{(d)})$.
$U'^{(d)}$ is a dense open subset of $C'^{(d)}$, hence $C'^{(d)}$ is a normal proper model of $U'^{(d)}$.
Let $C'^{[d]}$ be the normalization of $C^{(d)}$ in the function field $K(U'^{[d]})$,
then $q_U$ extends to a map $q: C'^{(d)} \to C'^{[d]}$, and we have the following diagrams:
\[
\begin{tikzcd}
	U' \arrow[r, hook] \arrow[d, "\pi_U"'] & C' \arrow[d, "\pi"] \\
	U \arrow[r, hook] & C
\end{tikzcd}
\qquad \qquad
\begin{tikzcd}
	{U'^{(d)}} \arrow[r, hook] \arrow[d, "q_U"'] & {C'^{(d)}} \arrow[d, "q"] \\
	{U'^{[d]}} \arrow[r, hook] \arrow[d] & {C'^{[d]}} \arrow[d] \\
	{U^{(d)}} \arrow[r, hook] & {C^{(d)}}
\end{tikzcd}
\]

\begin{lemma}\label{lemma:torsor_unramified_codim_one:reorg}
In the context of the preceding discussion, the morphism $U'^{(d)} \to U'^{[d]} \hookrightarrow C'^{[d]}$
is unramified in codimension $1$ (see \Cref{definition:tamely_ramified_in_codim_1:reorg}).
\end{lemma}
\begin{proof}

Let $C' \setminus U' = \{x_1, \ldots, x_m \}$. For each $x \in C' \setminus U'$, define the locus 
$$Z_x = \{D \in C'^{(d)} \mid x \in \mathrm{Supp}(D) \}.$$
Each $Z_x$ is a prime divisor in the symmetric product $C'^{(d)}$, and the boundary $C'^{(d)} \setminus U'^{(d)}$ is precisely the finite union $\bigcup_{i=1}^m Z_{x_i}$. 

Let $\eta$ be the generic point of $Z_x$. 
It suffices to show that the local ring extension $\mathcal{O}_{C'^{[d]}, q(\eta)} \hookrightarrow \mathcal{O}_{C'^{(d)}, \eta}$ is weakly unramified (\Cref{definition:weakly_unramified_and_tamely_ramified_DVR:reorg}).

Let $p: C'^d \to C'^{(d)}$ be the natural quotient. Consider the closed subvariety
$$H = \{x\} \times C' \times \cdots \times C' \subset C'^d.$$

The subvariety $H$ is irreducible, and its image under $p$ is exactly $Z_x$. 
Consequently, $p$ maps the generic point $\eta_H$ of $H$ to the generic point $\eta$ of $Z_x$.
\[
\begin{tikzcd}
\eta_H \arrow[d] \arrow[r] & H \arrow[r] \arrow[d] & C'^d \arrow[d, "p"] \\
\eta \arrow[r] \arrow[d] & Z_x \arrow[r] \arrow[d] & C^{(d)} \arrow[d, "q"] \\
q(\eta) \arrow[r] & q(Z_x) \arrow[r] & C'^{[d]}
\end{tikzcd}
\]
It suffices to show that the local ring extension $\mathcal{O}_{C'^{[d]}, q(\eta)} \hookrightarrow \mathcal{O}_{C'^d, \eta_H}$ is weakly unramified.
Let $R = \mathcal{O}_{C'^d, \eta_H}$. And let $A=\mathcal{O}_{C'^{[d]}, q(\eta)}$.

Recall that $C'^{[d]}$ is the quotient of $C'^d$ by the finite group $\Xi = \KG \rtimes S_d$ (see \Cref{cor:sym-power-quotient}),
and the morphism $q \circ p : C'^d \to C'^{[d]}$ is \textit{branched Galois Cover} (finite, surjective map of normal integral schemes $X \to Y$, with function field extension $K(X)/K(Y)$ Galois).

Hence the decomoposition group of $\eta_H$ in $\Xi$ and the inertia group of $\eta_H$ in $\Xi$ are defined:
$$D_{\eta_H} = \{ g \in \Xi \mid g(\eta_H) = \eta_H \}$$
$$I_{\eta_H} = \{ g \in D_{\eta_H} \mid g^*(u) \equiv u \pmod{\mathfrak{m}_{R}} \text{ for all } u \in R \}$$

%Then $A=R^{\Xi}$ where $\Xi = \KG \rtimes S_d$\footnote{Recall $\KG = \{ (g_1, \dots, g_d) \in G^d \mid g_1 g_2 \dots g_d = 1_G \}$} is the finite group acting on $C'^d$ as in \Cref{cor:sym-power-quotient}.
\textbf{Claim: } The inertia group $I_{\eta_H}$ is trivial.

Let $K=K(C')$ be the function field of $C'$, then 
$K(C'^d) = K^{\otimes_k d} \cong Frac(K_1 \otimes_k K_2 \otimes_k \cdots \otimes_k K_d)$
Where $K_i$ is the copy of $K$ corresponding to the $i$-th factor of $C'^d$.
The closed subscheme $H$ corresponds to the prime ideal $\pp_H \subset \Oo_{C'^d}$, and $R=\Oo_{C'^d, \pp_H}$. 
The residue field of $R$ is exactly $K(H) \cong Frac(k(x) \otimes_k K_2 \otimes_k \cdots \otimes_k K_d)$ where $k(x)$ is the residue field of $x$.

The action of $\gamma = (g_1, \ldots, g_d, \sigma) \in \Xi$, via the pullback $(\gamma^*)$ on $K(C'^d)$ does two things (in that order):
\begin{enumerate}
    \item It permutes the tensor factors according to $\sigma$. A function strictly in $K_i$ is mapped to a function in $K_{\sigma(i)}$.
    \item It applies the field automorphism $g_i^* \in \mathrm{Aut}(K/k)$ to the respective factors.
\end{enumerate}
Let $\gamma \in I_{\eta_H} \subset \Xi$, Then the induced automorphism $\bar{\gamma}^*$ on $K(H)$ must be the identity map.
This implies that $\sigma$ must fix $i$ for every $i \geq 2$ (The $K_i$ are linearly independent), and hence $\sigma=id$.
The trivial action of $\bar{\gamma}^*$ restricts to a trivial action of $g_i^*$ on $K_{i}$ for every $i \geq 2$, and hence 
$g_i = 1$ for every $i \geq 2$ (since the action of $G$ on $K(C')$ is faithful).
Finally, since $g_1 g_2 \cdots g_d = 1$, we also have $g_1 = 1$.
Since the inertia group $I_{\eta_H}$ is trivial, the local ring extension $A \hookrightarrow R$ is weakly unramified.
(See for example \stackstag{09E3}, \stackstag{0BTF})

\end{proof}


\subsection{Cyclic extensions: Kummer, Artin--Schreier, and Artin--Schreier--Witt}
The three classes of cyclic extensions we use share one shape. In each, a
cyclic group $A$ sits in a short exact sequence of commutative group schemes
$0 \to A \to \cG \xrightarrow{\,\psi\,} \cG \to 0$ for the \'etale topology (a
Lang-type isogeny), so that $\psi$-torsors over a field $K$ are classified by
$\cG(K)/\psi\,\cG(K)$ via additive Hilbert~90; and the ramification of the
associated extension of a local field is read off from the pole order(s) of a
normalized defining datum. The following dictionary records the
correspondence; the three columns are quoted in detail below.

\begin{table}[h]
\centering\footnotesize
\renewcommand{\arraystretch}{1.35}
\begin{tabular}{@{}l>{\raggedright\arraybackslash}p{3.1cm}>{\raggedright\arraybackslash}p{3.1cm}>{\raggedright\arraybackslash}p{3.7cm}@{}}
\hline
 & \textbf{Kummer} & \textbf{Artin--Schreier} & \textbf{Artin--Schreier--Witt} \\
\hline
group $A$ & $\Z/e\Z$, $\ p \nmid e$ & $\Z/p\Z$ & $\Z/p^r\Z$ \\
on $K$ & contains $\mu_e$ & $\operatorname{char} K = p$ & $\operatorname{char} K = p$ \\
isogeny $\psi$ & $t \mapsto t^{e}$ on $\bG_m$ & $\wp = t^{p}-t$ on $\bG_a$ & $\wp = F - \mathrm{id}$ on $\bW_r$ \\
sequence & $1 \to \mu_e \to \bG_m \xrightarrow{e} \bG_m \to 1$ & $0 \to \Z/p\Z \to \bG_a \xrightarrow{\wp} \bG_a \to 0$ & $0 \to \Z/p^r\Z \to \bW_r \xrightarrow{\wp} \bW_r \to 0$ \\
classifies & $K^\times/(K^\times)^e$ & $K/\wp K$ & $W_r(K)/\wp W_r(K)$ \\
break from & $v_K(a) \bmod e$ & pole order $m = -v_K(a)$ & Schmid--Witt $\max\limits_j p^{\,r-1-j} m_j$ \\
reference & \Cref{theorem:kummer_ramification:reorg} & \Cref{theorem:artin_schreier_ramification:reorg} & \Cref{theorem:asw_theory:reorg}, \Cref{theorem:schmid_witt_break_formula:reorg} \\
\hline
\end{tabular}
\caption{The three cyclic-extension theories in parallel.}
\label{table:cyclic_theory_dictionary:reorg}
\end{table}

\subsubsection*{Kummer and Artin--Schreier theory}
For cyclic extensions of degree prime to $p$ (Kummer) and of degree $p$
(Artin--Schreier), the following two theorems make the ``break from'' row of
\Cref{table:cyclic_theory_dictionary:reorg} precise: the valuation of a
defining element $a \in K$ determines the ramification index and the jumps in
the upper numbering filtration.
Throughout this section, let $K$ be a discrete valuation field with a perfect residue field $\kappa$ of characteristic $p > 0$.

\begin{theorem}[Ramification in Kummer Extensions, \cite{koch1997algebraic}, Proposition 1.83]\label{theorem:kummer_ramification:reorg}
%Let $K$ be a local field of characteristic $\text{char}(K) =p$  
Assume $K$ contains the $n$-th roots of unity $\mu_n$. 
Let $L/K$ be the extension given by the equation $X^n = a$ for some 
$a \in K^\times$ and denote by $G$ its Galois group. Then we have:

\begin{enumerate}
    \item If $v_K(a) \in n\mathbb{Z}$ and the image of $a\pi^{-v_K(a)}$ in the residue field $\kappa$ is an $n$-th power, the extension $L/K$ is trivial.
    \item If $v_K(a) \in n\mathbb{Z}$ and the image of $a\pi^{-v_K(a)}$ in the residue field $\kappa$ is not an $n$-th power, the extension $L/K$ is cyclic and unramified.
    \item If $v_K(a) \notin n\mathbb{Z}$, the extension $L/K$ is cyclic and ramified. Specifically, if $\gcd(|v_K(a)|, n) = 1$, the extension is totally ramified of degree $n$.
     Otherwise it has ramification index $\frac{n}{|gcd(v_K(a)|, n)}$
\end{enumerate}
Conversly, Kummer theory ensures that every cyclic extension of degree $n$, prime to $p$ of a field that contains $n$-th roots of unity, is of the above form.
Moreover, in the above we can always take $a \in \Oo_K$
\end{theorem}
Note that in the case of total ramification the extesntion is tamely ramified.


\begin{theorem}[Ramification in Artin-Schreier Extensions, \cite{Thomas2005}]\label{theorem:artin_schreier_ramification:reorg}
Let $\wp(x) = x^p - x$ be the Artin-Schreier opeartor.
Let $L/K$ be the extension given by the equation $X^p - X = a$ for some $a \in K$ and denote by $G$ its Galois group. 
Then we have:
\begin{enumerate}
    \item If $v_K(a) > 0$ or if $v_K(a) = 0$ and $a \in \wp(K)$, the extension $L/K$ is trivial.
    \item If $v_K(a) = 0$ and if $a \notin \wp(K)$, the extension $L/K$ is cyclic of degree $p$ and unramified.
    \item If $v_K(a) = -m < 0$ with $m \in \mathbb{Z}_{>0}$ and if $m$ is prime to $p$, the extension $L/K$ is cyclic of degree $p$ again and totally ramified. Moreover,
     its ramification groups are given by:
     $$G = G^{(-1)} = \dots = G^{(m)} \quad \text{and} \quad G^{(m+1)} = 1.$$
\end{enumerate}
Conversely, Artin-Schreier theory ensures that every cyclic extension of degree $p$ takes this form.
Moreover, in the above and under the isomorphism $K \cong k((t))$, one can always take $a$ of the form 
$c t^{-m} + a_{-m + 1} t^{-m + 1} + \cdots a_{-1}t^{-1} + a_0$. If $k$ is algebraically closed then there is a 
change of variables such that $a=u^{-m}$.
\end{theorem}


% =====================================================================
%  Witt-vector prerequisites
%  ---------------------------------------------------------------
%  This file collects the basic theory of ($p$-typical, finite-length)
%  Witt vectors and of Artin--Schreier--Witt theory used in the proof
%  that $\cP^{[\deg \ndls]}$ is unramified at $\eta_\ndls$ for
%  $G = \Z/p^r\Z$.  It is meant to be \input inside the Preliminaries.
%
%  Macros: if \bW is not already defined in MacrosAndFigures.tex,
%  the \providecommand below supplies it (consider moving it there).
% =====================================================================
\providecommand{\bW}{\mathbb{W}}

\subsubsection*{Artin--Schreier--Witt theory}
\phantomsection\label{subsection:witt_vector_prerequisites:reorg}
This is the third column of \Cref{table:cyclic_theory_dictionary:reorg}: it
generalizes Artin--Schreier theory from $\Z/p\Z$ to $\Z/p^r\Z$ by replacing the
field $K$ with the ring $W_r(K)$ of length-$r$ Witt vectors, the operator
$\wp(u)=u^p-u$ with $\wp = F-\mathrm{id}$, and the pole-order break criterion of
\Cref{theorem:artin_schreier_ramification:reorg} with the Schmid--Witt formula
(\Cref{theorem:schmid_witt_break_formula:reorg}). The development is
correspondingly longer, since the ring $W_r$ itself must be set up; a reader
willing to take \Cref{theorem:asw_theory:reorg} and
\Cref{theorem:schmid_witt_break_formula:reorg} on faith may skip to
\Cref{theorem:ramification_after_blowup_P1:reorg}.

Artin--Schreier theory, together with the ramification formula of
\Cref{theorem:artin_schreier_ramification:reorg}, gives complete control over the cyclic
extensions of degree $p$ of a local field of characteristic $p$. To treat the groups
$G = \Z/p^r\Z$ we need the analogous control over cyclic extensions of degree $p^r$,
which is provided by Witt's generalization: the field $K$ is replaced by the ring
$W_r(K)$ of $p$-typical Witt vectors of length $r$, the operator $\wp(u) = u^p - u$ by
its Witt-vector analogue $\wp = F - \mathrm{id}$, and the classical ramification-break
formula by the Schmid--Witt formula. This subsection collects the definitions and the
elementary lemmas that we use; our references for the general theory are
\cite[II, \S6]{serreLF} and \cite{Rabinoff2014Witt}.

Throughout, $p$ is a fixed prime, and $r \geq 1$ is an integer. All the Witt vectors
appearing in this thesis are $p$-typical and of finite length.

\subsubsection*{The ring of Witt vectors}

For $n \geq 0$, the $n$-th \textbf{Witt polynomial} in the variables $X_0, \dots, X_n$
is
\[
w_n(X_0, \dots, X_n) \;=\; \sum_{j=0}^{n} p^j X_j^{p^{n-j}}
\;=\; X_0^{p^n} + p X_1^{p^{n-1}} + \dots + p^n X_n .
\]

\begin{theorem}[Witt; {\cite[II, \S6, Theorem~6]{serreLF}}, {\cite[\S1]{Rabinoff2014Witt}}]
\label{theorem:witt_ring_structure:reorg}
There exists a unique family of polynomials
$S_n, P_n \in \Z[X_0, \dots, X_n; Y_0, \dots, Y_n]$, $n \geq 0$, such that
\[
w_n(S_0, \dots, S_n) = w_n(X_0,\dots,X_n) + w_n(Y_0,\dots,Y_n),
\qquad
w_n(P_0, \dots, P_n) = w_n(X_0,\dots,X_n) \cdot w_n(Y_0,\dots,Y_n)
\]
for all $n$. For any commutative ring $A$, the operations
\[
\vec a \oplus \vec b := \bigl(S_n(\vec a; \vec b)\bigr)_{0 \le n \le r-1},
\qquad
\vec a \odot \vec b := \bigl(P_n(\vec a; \vec b)\bigr)_{0 \le n \le r-1}
\]
make the set $A^r$ into a commutative ring $W_r(A)$, the ring of \textbf{$p$-typical
Witt vectors of length $r$ over $A$}, with zero $(0, \dots, 0)$ and unit
$(1, 0, \dots, 0)$. Additive inverses are likewise given by universal integral
polynomials; we write $\ominus$ for Witt subtraction. The \textbf{ghost map}
\[
w \colon W_r(A) \longrightarrow A^r,
\qquad
\vec a \longmapsto \bigl(w_0(\vec a), w_1(\vec a), \dots, w_{r-1}(\vec a)\bigr)
\]
is a ring homomorphism to the product ring $A^r$; it is injective when $p$ is a
non-zero-divisor in $A$, and bijective when $p \in A^\times$.
\end{theorem}

We index the components of a Witt vector $\vec a = (a_0, \dots, a_{r-1})$ by
$0, \dots, r-1$. For $r = 1$ the polynomials reduce to $S_0 = X_0 + Y_0$ and
$P_0 = X_0 Y_0$, so $W_1(A) = A$ as a ring.

Since the ring operations are given by universal polynomials with integer
coefficients, $W_r$ is a functor: a ring homomorphism $\sigma \colon A \to B$ induces
the ring homomorphism
\[
W_r(\sigma) \colon W_r(A) \to W_r(B),
\qquad
W_r(\sigma)(a_0, \dots, a_{r-1}) = (\sigma a_0, \dots, \sigma a_{r-1}),
\]
acting \emph{componentwise}. Two consequences of this trivial-looking observation are
used repeatedly and deserve emphasis: a ring automorphism $\sigma$ of $A$ commutes
with $\oplus$ and $\ominus$ on $W_r(A)$, and a Witt vector $\vec a \in W_r(A)$
satisfies $W_r(\sigma)(\vec a) = \vec a$ if and only if \emph{each component} $a_n$ is
$\sigma$-invariant. In particular, for a group $\Gamma$ acting on $A$ by ring
automorphisms,
\begin{equation}\label{eq:witt_invariants_componentwise:reorg}
W_r(A)^{\Gamma} \;=\; W_r\bigl(A^{\Gamma}\bigr)
\qquad \text{inside } W_r(A).
\end{equation}

The uniqueness statement in \Cref{theorem:witt_ring_structure:reorg} is an instance of a
general principle that we invoke several times, so we record it explicitly.

\begin{lemma}[ghost principle]\label{lemma:ghost_principle:reorg}
Let $\Phi$ and $\Psi$ be two operations on Witt vectors given by universal polynomials
with integer coefficients in the components of their arguments. If
$w \circ \Phi = w \circ \Psi$ holds over the polynomial ring
$\Z[X_0, X_1, \dots]$ on the components of the arguments, then $\Phi = \Psi$ over
every commutative ring.
\end{lemma}

\begin{proof}
Over $\Z[X_0, X_1, \dots]$ the prime $p$ is a non-zero-divisor, so the ghost map is
injective by \Cref{theorem:witt_ring_structure:reorg} and $\Phi = \Psi$ there; this is an
identity of integral polynomials, which persists under every base change
$\Z[X_0, X_1, \dots] \to A$.
\end{proof}

The only structural feature of the addition polynomials $S_n$ that we shall need is
their ``triangular'' shape:

\begin{prop}[shape of Witt addition]\label{prop:witt_addition_shape:reorg}
For every $n \geq 0$ there is a polynomial
$C_n \in \Z[X_0, \dots, X_{n-1}; Y_0, \dots, Y_{n-1}]$, involving only the variables of
index \emph{strictly less} than $n$, such that
\begin{equation}\label{eq:witt_addition_shape:reorg}
(\vec a \oplus \vec b)_n \;=\; a_n + b_n + C_n(a_0, \dots, a_{n-1}; b_0, \dots, b_{n-1}).
\end{equation}
We refer to $C_n$ as the \textbf{carry polynomial}. The analogous statement holds for
$\ominus$.
\end{prop}

\begin{proof}
By induction on $n$, using the recursion defining $S_n$:
\[
p^n S_n \;=\; \sum_{j \leq n} p^j \bigl(X_j^{p^{n-j}} + Y_j^{p^{n-j}}\bigr)
\;-\; \sum_{j < n} p^j S_j^{p^{n-j}} .
\]
The terms with $j = n$ contribute $p^n (X_n + Y_n)$, and all remaining terms involve
only variables and polynomials $S_j$ of index $j < n$, which by induction lie in
$\Z[X_{<n}; Y_{<n}]$. Since $S_n$ has integer coefficients
(\Cref{theorem:witt_ring_structure:reorg}), $C_n := S_n - X_n - Y_n$ is an integral
polynomial in the variables of index $< n$. The statement for $\ominus$ follows by the
same argument applied to the inversion polynomials.
\end{proof}

\subsubsection*{Frobenius, Verschiebung and truncation}

From now on, $A$ denotes an $\F_p$-algebra; this is the only case we use. Three
operators act on the rings $W_r(A)$:

\begin{definition}\label{definition:witt_operators:reorg}
Let $A$ be an $\F_p$-algebra.
\begin{enumerate}
    \item The \textbf{Frobenius} $F \colon W_r(A) \to W_r(A)$ is
    $F := W_r(\varphi_A)$, where $\varphi_A \colon a \mapsto a^p$ is the absolute
    Frobenius of $A$. Explicitly, $F$ acts componentwise:
    $F(a_0, \dots, a_{r-1}) = (a_0^p, \dots, a_{r-1}^p)$. By functoriality, $F$ is a
    ring endomorphism of $W_r(A)$.
    \item The \textbf{Verschiebung} is the shift
    \[
    V \colon W_{r-1}(A) \to W_r(A),
    \qquad
    V(a_0, \dots, a_{r-2}) = (0, a_0, \dots, a_{r-2}).
    \]
    Iterating, $V^j \colon W_{r-j}(A) \to W_r(A)$; in particular
    $V^{r-1} \colon A = W_1(A) \to W_r(A)$ sends $c \mapsto (0, \dots, 0, c)$.
    \item The \textbf{truncation} is
    \[
    t \colon W_r(A) \to W_{r-1}(A),
    \qquad
    t(a_0, \dots, a_{r-1}) = (a_0, \dots, a_{r-2}).
    \]
\end{enumerate}
\end{definition}

\begin{rem*}
For a general (not necessarily $\F_p$-) algebra $A$, the Frobenius is defined by
universal polynomials characterized by $w_n \circ F = w_{n+1}$, and the componentwise
formula above holds precisely in characteristic $p$; see
\cite[II, \S6]{serreLF} or \cite[\S2]{Rabinoff2014Witt}. We will not need the general
case.
\end{rem*}

We record from the general theory (\cite[II, \S6]{serreLF},
\cite[\S2]{Rabinoff2014Witt}) how multiplication by $p$ interacts with these
operators.

\begin{prop}\label{prop:witt_multiplication_by_p:reorg}
Let $A$ be an $\F_p$-algebra. Then:
\begin{enumerate}
    \item $t$ is a natural ring homomorphism commuting with $F$, and its kernel is
    \[
    \ker\bigl(t \colon W_r(A) \to W_{r-1}(A)\bigr)
    \;=\; V^{r-1} A \;=\; \{(0, \dots, 0, c) : c \in A\}.
    \]
    \item $V$ is additive, and $F V = V F$.
    \item Multiplication by $p$ on $W_r(A)$ equals $V \circ F = F \circ V$
    (composed with the appropriate truncation); explicitly,
    \[
    p \cdot (a_0, \dots, a_{r-1}) \;=\; (0,\, a_0^p,\, \dots,\, a_{r-2}^p).
    \]
    Iterating, $p^j \cdot \vec a = (0, \dots, 0, a_0^{p^j}, \dots, a_{r-1-j}^{p^j})$;
    in particular $p^r = 0$ on $W_r(A)$, and
    \begin{equation}\label{eq:witt_p_to_r_minus_1:reorg}
    p^{r-1} \cdot (a_0, \dots, a_{r-1}) \;=\; \bigl(0, \dots, 0,\, a_0^{p^{r-1}}\bigr).
    \end{equation}
\end{enumerate}
\end{prop}

\begin{proof}
(1) That $t$ is a ring homomorphism follows from the recursive definition of the
$S_n, P_n$: the first $r-1$ components of a sum or product depend only on the first
$r-1$ components of the arguments (\Cref{prop:witt_addition_shape:reorg}). Naturality and
the commutation with the (componentwise) $F$ are clear, and the kernel is computed
componentwise. (2) Additivity of $V$ follows from \Cref{lemma:ghost_principle:reorg} and
the ghost identity $w_n(V\vec a) = p\, w_{n-1}(\vec a)$ (with $w_{-1} := 0$); the
commutation $FV = VF$ is immediate since $F$ is componentwise. (3) The identity
$F V = p$ holds over arbitrary rings for the general Frobenius, by
\Cref{lemma:ghost_principle:reorg} applied to
$w_n(FV\vec a) = w_{n+1}(V \vec a) = p\, w_n(\vec a) = w_n(p \vec a)$; in
characteristic $p$ the general Frobenius is componentwise, giving the displayed
formula. See \cite[II, \S6, Proposition~8 ff.]{serreLF} for details.
\end{proof}

\begin{example}\label{example:witt_of_Fp:reorg}
$W_r(\F_p) \cong \Z/p^r\Z$ canonically. Indeed, by
\Cref{eq:witt_p_to_r_minus_1:reorg} the unit $\underline 1 = (1, 0, \dots, 0)$ satisfies
$p^j \cdot \underline 1 = (0, \dots, 0, 1, 0, \dots, 0)$ ($1$ in position $j$), so
$\underline 1$ has additive order $p^r$; since $|W_r(\F_p)| = p^r$, the additive group
is cyclic of order $p^r$ generated by the unit, and $m \mapsto m \cdot \underline 1$
is a ring isomorphism $\Z/p^r\Z \xrightarrow{\;\cong\;} W_r(\F_p)$. We write
$\underline m \in W_r(\F_p)$ for the image of $m \in \Z/p^r\Z$. (Similarly, the full
ring of Witt vectors satisfies $W(\F_p) \cong \Z_p$; we shall not need this.)
\end{example}

\begin{definition}\label{definition:asw_operator:reorg}
The \textbf{Artin--Schreier--Witt operator} on $W_r(A)$, $A$ an $\F_p$-algebra, is
\[
\wp \;:=\; F - \mathrm{id},
\qquad
\wp\,\vec u \;=\; F\vec u \ominus \vec u .
\]
For $r = 1$ this is the classical operator $\wp(u) = u^p - u$ of
\Cref{theorem:artin_schreier_ramification:reorg}.
\end{definition}

Since $F$ is a ring endomorphism, $\wp$ is an endomorphism of the additive group
$(W_r(A), \oplus)$:
$\wp(\vec u \oplus \vec v) = \wp\vec u \oplus \wp\vec v$. Moreover $\wp$ commutes with
$W_r(\sigma)$ for every ring homomorphism $\sigma$ (both $F$ and the ring operations
are natural), and with $V$ and $t$ by \Cref{prop:witt_multiplication_by_p:reorg}.

The following bookkeeping lemma is elementary but does the essential work of reducing
statements about $W_r$ to statements about $W_{r-1}$ and about $A$ itself: the
truncation $t$ and the inclusion $V^{r-1}$ of its kernel slice $W_r$ into lower-length
pieces, compatibly with $\wp$, and \emph{addition of an element of the kernel is
carry-free}.

\begin{lemma}[Witt bookkeeping]\label{lemma:witt_bookkeeping:reorg}
Let $A$ be an $\F_p$-algebra.
\begin{enumerate}[label=(\roman*)]
    \item $t \colon W_r(A) \to W_{r-1}(A)$ is a ring homomorphism commuting with $F$,
    hence with $\wp$; its kernel is $V^{r-1} A$.
    \item $V^{r-1} \colon A \to W_r(A)$ is additive, and
    $\wp \circ V^{r-1} = V^{r-1} \circ \wp$, where on the left $\wp$ is the operator
    on $W_r(A)$ and on the right the classical $\wp(c) = c^p - c$ on $A$.
    \item (carry-free top-level addition) For all $\vec a \in W_r(A)$ and $c \in A$,
    \[
    \vec a \oplus V^{r-1}(c) \;=\; (a_0, \dots, a_{r-2},\, a_{r-1} + c).
    \]
\end{enumerate}
\end{lemma}

\begin{proof}
(i) and (ii) restate parts of \Cref{prop:witt_multiplication_by_p:reorg} (for (ii): $V^{r-1}$
is additive and commutes with the componentwise $F$, hence with
$\wp = F - \mathrm{id}$).
(iii) By \Cref{lemma:ghost_principle:reorg} it suffices to verify the identity over
$\Z[a_0, \dots, a_{r-1}, c]$, where the ghost map is injective. Put
$\vec b := (0, \dots, 0, c)$ and $\vec s := \vec a \oplus \vec b$. For $n < r-1$ we
have $w_n(\vec b) = 0$, so $w_n(\vec s) = w_n(\vec a)$, and by induction on $n$
(solving for $s_n$, using that $p$ is a non-zero-divisor) $s_n = a_n$ for
$n \leq r-2$. For $n = r-1$, $w_{r-1}(\vec b) = p^{r-1} c$ gives
$p^{r-1} s_{r-1} = p^{r-1} a_{r-1} + p^{r-1} c$, i.e.\ $s_{r-1} = a_{r-1} + c$.
\end{proof}

\begin{rem*}
\Cref{lemma:witt_bookkeeping:reorg}(iii) fails for the addition of vectors supported in
lower positions: in general the carry polynomials $C_n$ of
\Cref{eq:witt_addition_shape:reorg} are non-zero, and Witt addition is \emph{not}
componentwise. In particular, operations performed componentwise on a Witt vector ---
such as discarding the regular parts of Laurent-series components --- do \emph{not}
in general respect Witt addition, and will be carefully avoided.
\end{rem*}

\subsubsection*{The Artin--Schreier--Witt isogeny and its torsors}

The componentwise ring operations make $\A^r_{\F_p}$ a ring scheme, the
\textbf{Witt vector group scheme} $\bW_r$: its underlying scheme is
$\A^r_{\F_p} = \Spec \F_p[T_0, \dots, T_{r-1}]$, with the group law given by the
addition polynomials $S_n$, so that $\bW_r(A) = W_r(A)$ functorially. The operator
$\wp = F - \mathrm{id}$ is an endomorphism of the group scheme $\bW_r$, and it is the
Witt-vector instance of a \emph{Lang isogeny}: by \Cref{example:witt_of_Fp:reorg} and the
lemma below, it sits in a short exact sequence of group schemes over $\F_p$ for the
\'etale topology
\begin{equation}\label{eq:asw_lang_sequence:reorg}
0 \longrightarrow \Z/p^r\Z \cong W_r(\F_p)
\longrightarrow \bW_r \xrightarrow{\;\wp\;} \bW_r \longrightarrow 0 .
\end{equation}

The corresponding covers admit the following completely explicit description, which
is the form in which we use them. For $\vec g \in W_r(A)$ we write
$A[\vec X]/(\wp \vec X \ominus \vec g)$, $\vec X = (X_0, \dots, X_{r-1})$, for the
quotient of $A[X_0, \dots, X_{r-1}]$ by the ideal generated by the $r$ components of
the Witt vector $\wp \vec X \ominus \vec g \in W_r(A[X_0, \dots, X_{r-1}])$.

\begin{lemma}[ASW covers are \'etale torsors]\label{lemma:asw_etale_torsor:reorg}
Let $A$ be an $\F_p$-algebra and $\vec g \in W_r(A)$. Let
\[
B \;:=\; A[\vec X]\big/\bigl(\wp \vec X \ominus \vec g\bigr).
\]
Then $B$ is finite free over $A$ with basis
$\bigl\{\prod_{n} X_n^{a_n} : 0 \le a_n \le p-1\bigr\}$, and \'etale over $A$. The
translation action of $W_r(\F_p) \cong \Z/p^r\Z$,
\[
m \colon \vec X \longmapsto \vec X \oplus \underline m ,
\]
makes $\Spec B$ a $\Z/p^r\Z$-torsor over $\Spec A$.
\end{lemma}

\begin{proof}
By \Cref{eq:witt_addition_shape:reorg} and the componentwise formula for $F$, the $n$-th
defining equation has the triangular form
\[
X_n^p - X_n - g_n + C_n'(X_0, \dots, X_{n-1};\, g_0, \dots, g_{n-1}) \;=\; 0
\]
for universal integral polynomials $C_n'$. Triangularity gives the free basis; the
Jacobian matrix of the system is lower triangular with diagonal entries
$p X_n^{p-1} - 1 = -1$, a unit, so $B$ is \'etale over $A$. Finally, $\Spec B$ is the
pullback along $\vec g \colon \Spec A \to \bW_r$ of the covering
$\wp \colon \bW_r \to \bW_r$; by the triangular form just established (over
$A = \F_p[T_0,\dots,T_{r-1}]$, $\vec g = (T_0,\dots,T_{r-1})$), $\wp$ is a surjective
\'etale homomorphism of group schemes, and its kernel is
$\{\vec u : F \vec u = \vec u\} = W_r(\F_p)$ since $F$ is componentwise. A surjective
\'etale group homomorphism is a torsor under its kernel, and torsors pull back to
torsors; hence $\Spec B$ is a $W_r(\F_p)$-torsor over $\Spec A$.
\end{proof}

\subsubsection*{Artin--Schreier--Witt theory for fields}

Over a field, the sequence \Cref{eq:asw_lang_sequence:reorg} computes all cyclic extensions
of $p$-power degree, exactly as classical Artin--Schreier theory does for degree $p$.

\begin{theorem}[Artin--Schreier--Witt theory; {\cite[\S3--4]{Thomas2005}}]
\label{theorem:asw_theory:reorg}
Let $M \supseteq \F_p$ be a field, $M^{s}$ a separable closure, and
$\Gamma = \Gal(M^{s}/M)$. Then $\wp$ is surjective on $W_r(M^{s})$ with kernel
$W_r(\F_p)$, and there is a canonical isomorphism of abelian groups
\[
\delta \colon\; W_r(M)\big/\wp\, W_r(M) \;\xrightarrow{\;\cong\;}\;
\Hom_{\mathrm{cont}}\bigl(\Gamma,\, \Z/p^r\Z\bigr) \;=\; H^1\bigl(M, \Z/p^r\Z\bigr),
\qquad
[\vec g] \longmapsto \bigl(\sigma \mapsto \sigma(\vec u) \ominus \vec u\bigr),
\]
where $\vec u \in W_r(M^{s})$ is any solution of $\wp \vec u = \vec g$. Under the
correspondence between characters and torsors
(\Cref{cor:abelian_equivilance_fundamental_torsors}), the class $[\vec g]$ corresponds
to the $\Z/p^r\Z$-torsor
\[
\Spec\, M[\vec X]\big/\bigl(\wp \vec X \ominus \vec g\bigr)
\]
of \Cref{lemma:asw_etale_torsor:reorg}. Moreover, for a class of order $p^s$ in
$W_r(M)/\wp W_r(M)$, the fixed field $M(\vec u) := M(u_0, \dots, u_{r-1})$ of
$\ker \delta[\vec g]$ is a cyclic extension of $M$ of degree $p^s$; the torsor is
connected if and only if $[\vec g]$ has order $p^r$, in which case
$\Spec M[\vec X]/(\wp \vec X \ominus \vec g) \cong \Spec M(\vec u)$ is (the spectrum
of) a cyclic field extension of degree $p^r$. Every cyclic extension of $M$ of degree
dividing $p^r$ arises this way.
\end{theorem}

\begin{proof}[Sketch of proof]
Surjectivity of $\wp$ on $W_r(M^{s})$ and the computation of its kernel follow from
the triangular equations in the proof of \Cref{lemma:asw_etale_torsor:reorg}: each component
of a preimage is obtained by solving a separable equation $X_n^p - X_n = (\dots)$, and
$F\vec u = \vec u$ holds if and only if $u_n^p = u_n$ for all $n$, i.e.\
$\vec u \in W_r(\F_p)$. Given $\vec g$ and a solution $\wp\vec u = \vec g$, for
$\sigma \in \Gamma$ we have
$\wp(\sigma \vec u \ominus \vec u) = \sigma \vec g \ominus \vec g = 0$ (using
\Cref{eq:witt_invariants_componentwise:reorg} and that $W_r(\sigma)$ commutes with $\wp$),
so $\sigma \vec u \ominus \vec u \in W_r(\F_p)$; one checks it is a continuous
homomorphism in $\sigma$ (the kernel is fixed pointwise by $\Gamma$), independent of
the choices modulo nothing and of the representative $\vec g$ modulo $\wp W_r(M)$.
Injectivity: if $\delta[\vec g] = 0$ then $\vec u$ is $\Gamma$-invariant, hence lies
in $W_r(M)$ by \Cref{eq:witt_invariants_componentwise:reorg}, so
$\vec g = \wp\vec u \in \wp W_r(M)$. Surjectivity of $\delta$ is the additive
Hilbert~90: $H^1(\Gamma, W_r(M^{s})) = 0$, by induction on $r$ along the exact
sequence $0 \to M^{s} \xrightarrow{V^{r-1}} W_r(M^{s}) \xrightarrow{t} W_{r-1}(M^{s})
\to 0$ of $\Gamma$-modules (\Cref{lemma:witt_bookkeeping:reorg}(i)) and the classical case
$H^1(\Gamma, M^{s}) = 0$. The remaining statements follow from
\Cref{theorem:equivalence_fundamental_covers} applied to the character
$\delta[\vec g]$, whose image has the same order as the class $[\vec g]$. For details
see \cite[\S3--4]{Thomas2005} or \cite{Rabinoff2014Witt}.
\end{proof}

The order of a class is partially visible in its $0$-th component:

\begin{cor}\label{cor:asw_class_order_and_first_component:reorg}
Let $M \supseteq \F_p$ be a field and $\vec g \in W_r(M)$ with $g_0 = 0$. Then
$p^{r-1}[\vec g] = 0$ in $W_r(M)/\wp W_r(M)$; equivalently, the cyclic extension
attached to $[\vec g]$ has degree at most $p^{r-1}$. Consequently, if $[\vec g]$ has
order $p^r$ (e.g.\ if $\vec g$ defines a connected torsor), then \emph{every}
representative of the class has non-zero $0$-th component.
\end{cor}

\begin{proof}
By \Cref{eq:witt_p_to_r_minus_1:reorg}, $p^{r-1} \vec g = (0, \dots, 0, g_0^{p^{r-1}}) = 0$
already in $W_r(M)$.
\end{proof}

\begin{rem*}\label{rem:asw_local_totally_ramified:reorg}
When $M = k((x))$ with $k$ algebraically closed, every finite separable extension
$L/M$ is totally ramified: the residue field extension is trivial (as $k$ is
algebraically closed) and $\Oo_L$ is a complete discrete valuation ring, so
$[L : M] = ef = e$. Thus a class of order $p^r$ in $W_r(k((x)))/\wp$ corresponds to a
cyclic, \emph{totally ramified} extension of degree $p^r$.
\end{rem*}

\subsubsection*{Isobaricity of the Witt addition polynomials}

The following homogeneity property of the universal addition polynomials is what
makes degree estimates on Witt sums possible: even though the carry polynomials $C_n$
of \Cref{eq:witt_addition_shape:reorg} are complicated, they are \emph{isobaric} for the
weighting in which the $j$-th component has weight $p^j$.

\begin{lemma}[isobaricity of Witt addition]\label{lemma:witt_addition_isobaric:reorg}
Assign to the variables $X_j$ and $Y_j$ the weight $p^j$. Then the $n$-th addition
polynomial $S_n$ is isobaric of weight $p^n$: every monomial of $S_n$ has total weight
exactly $p^n$. The same holds for each component of a $d$-fold iterated Witt sum
$\bigoplus_{i=1}^{d} \vec a_i$, viewed as a universal polynomial in the components
$a_{i,j}$ of weight $p^j$.
\end{lemma}

\begin{proof}
The $n$-th ghost polynomial $w_n = \sum_{j \le n} p^j X_j^{p^{n-j}}$ is isobaric of
weight $p^n$ for this weighting. The recursion
\[
p^n S_n = \sum_{j \le n} p^j \bigl(X_j^{p^{n-j}} + Y_j^{p^{n-j}}\bigr)
- \sum_{j < n} p^j S_j^{p^{n-j}}
\]
preserves isobaricity of weight $p^n$ by induction on $n$: each $S_j$ is isobaric of
weight $p^j$, so $S_j^{p^{n-j}}$ is isobaric of weight $p^n$. Isobaricity persists
under iterated substitution of Witt sums into Witt sums (substituting an isobaric
polynomial of weight $p^j$ for a variable of weight $p^j$ preserves weights), giving
the statement for $d$-fold sums.
\end{proof}

\subsubsection*{Ramification breaks: the Schmid--Witt formula}

Finally we state the generalization of the ramification formula of
\Cref{theorem:artin_schreier_ramification:reorg} to Artin--Schreier--Witt extensions. As in
the classical case, the formula reads off the largest upper-numbering ramification
break from the pole orders of a defining Witt vector --- provided the representative
is normalized so that no interference between the levels can occur. The relevant
normalization is the following.

\begin{definition}\label{definition:reduced_witt_vector:reorg}
Let $K = k((x))$. For $g \in K$ put $m(g) := \max\bigl(0, -v_x(g)\bigr)$ (the pole
order of $g$, or $0$ if $g$ is regular). A vector
$\vec g = (g_0, \dots, g_{r-1}) \in W_r(K)$ is \textbf{reduced} if for every $j$,
either $m(g_j) = 0$ or $p \nmid m(g_j)$.
\end{definition}

\begin{theorem}[Schmid--Witt break formula]\label{theorem:schmid_witt_break_formula:reorg}
Let $K = k((x))$ with $k$ perfect of characteristic $p$, and let $L/K$ be a cyclic
totally ramified extension of degree $p^r$ whose Artin--Schreier--Witt class
(\Cref{theorem:asw_theory:reorg}) is represented by a \emph{reduced} vector
$\vec g \in W_r(K)$, with $m_j := m(g_j)$. Then the largest upper-numbering
ramification break of $L/K$ equals
\[
u \;=\; \max_{0 \le j \le r-1} \; p^{\,r-1-j}\, m_j .
\]
\end{theorem}

This is due, for finite residue fields, to Kanesaka--Sekiguchi
\cite{KanesakaSekiguchi1979} and Thomas \cite{Thomas2005} (building on Schmid
\cite{Schmid1937} for $r = 1$), and for arbitrary perfect residue fields to
Elder--Keating \cite{ElderKeating2025}. We use it as stated, in one direction only:
\emph{evaluating the formula on one reduced representative computes $u$}; no
minimization over representatives is needed.

\begin{rem*}
Reducedness is exactly what makes the formula exact. If two levels $j < j'$ with
$m_j, m_{j'} \neq 0$ satisfied $p^{r-1-j} m_j = p^{r-1-j'} m_{j'}$, then
$m_{j'} = p^{\,j'-j} m_j$ would be divisible by $p$, contradicting
\Cref{definition:reduced_witt_vector:reorg}; hence the maximum is attained at a unique
level, and no cancellation of breaks between levels can occur.
\end{rem*}

\begin{rem*}
For $r = 1$, a reduced vector is a single element $g \in K$ whose pole order $m_0$ is
zero or prime to $p$, and \Cref{theorem:schmid_witt_break_formula:reorg} recovers
\Cref{theorem:artin_schreier_ramification:reorg}: the largest (indeed the only) upper break
is $m_0$. The existence of reduced representatives --- and, over $K = k((x))$, of
reduced representatives whose components are polynomials in $x^{-1}$ without constant
term --- is not part of the general theory quoted here; it will be proved where it is
used, by induction on $r$ through the truncation maps of
\Cref{lemma:witt_bookkeeping:reorg}.
\end{rem*}


\section{Reduction to \texorpdfstring{$\bP^1_k$}{P1}}

We argue that it is enough to prove \Cref{theorem:torsors_after_blowup_is_tame:reorg} for $C=\bP^1_k$, $\ndls = d \cdot 0$, and $\cP \to \bG_m = \bP^1_k \setminus \{0, \infty\}$ a $G$-torsor given by an explicit global equation.
Informally, this is because \emph{the ramification of $\cP^{[d]}$ at $\eta_\ndls$ is a formal--local invariant of $\cP$ at $P$}:
it depends only on the completed local ring $\widehat{\Oo_{C,P}} \cong k[[x]]$ together with the restriction of $\cP$
to the punctured formal disc $\Spec(k((x)))$, since all of the constructions involved --- self products, quotients by finite
groups, and blowups --- commute with the flat base change to the completion; for blowups, this is
\Cref{theorem:blowup_flat_base_change}.
We emphasize that this is not a general principle that can simply be quoted: it is a statement about the particular
constructions at hand, and it is proved by checking each construction separately. This is the content of
\Cref{lemma:formal_local_invariance:reorg} and \Cref{cor:reduction_to_P1:reorg} below.
The second ingredient of the reduction is \Cref{lemma:realization_over_Gm:reorg}: every local torsor class appearing in
\Cref{theorem:torsors_after_blowup_is_tame:reorg} is realized by an explicit torsor over $\bG_m \subset \bP^1_k$.

Throughout this section we assume that $k$ is algebraically closed; in particular every closed point of $C$ is $k$-rational, and $\ndls = dP$ with $P \in C(k)$ and $d = \deg \ndls$.

\begin{rem*}
    The assumption that $k$ is algebraically closed is harmless in our situation, for the following reasons.
    \begin{enumerate}
        \item \textbf{(Base change to $\bar{k}$.)} Assume $k$ is perfect, so that $\bar{k}=k^{sep}$.
        The formation of $C^{(d)}$ (a quotient by a finite group), of the blowup $X_\ndls$ along $Z_\ndls$, and of
        $\cP^{[d]}$ (\Cref{cor:sym-power-quotient}) commutes with the flat base change $\Spec(\bar{k}) \to \Spec(k)$.
        Since $E_\ndls \cong \bP^{\deg \ndls - 1}_k$ is geometrically irreducible, $(E_\ndls)_{\bar{k}}$ is irreducible
        and its generic point $\bar{\eta}$ lies over $\eta_\ndls$; moreover
        $\Oo_{X_\ndls, \eta_\ndls} \to \Oo_{(X_\ndls)_{\bar{k}}, \bar{\eta}}$ is a weakly unramified extension of discrete
        valuation rings with separable residue field extension. The proof of \Cref{lemma:descend_ramification_along_etale:reorg}
        uses only these two properties of the extension (there the extension was moreover finite étale), so it applies
        verbatim: the ramification of $(\cP_{\bar{k}})^{[d]} = (\cP^{[d]})_{\bar{k}}$ at $\bar{\eta}$ is bounded by $r$ if
        and only if the ramification of $\cP^{[d]}$ at $\eta_\ndls$ is. Hence unramifiedness and tameness at $\eta_\ndls$
        may be checked after base change to $\bar{k}$.
        \item \textbf{(Rationality of $P$ --- a caveat.)} If $P$ is a closed point with $\kappa(P) \neq k$, then after base
        change to $\bar{k}$ the point $Z_{\ndls}$ corresponds to the divisor $\sum_{\bar{P} \mapsto P} n \bar{P}$ on
        $C_{\bar{k}}$, whose support consists of $[\kappa(P):k]$ distinct points, whereas the computation below treats a
        modulus supported at a single rational point. We therefore prove \Cref{theorem:torsors_after_blowup_is_tame:reorg} under
        the standing assumption that $k$ is algebraically closed; this suffices for our purposes, since in the proof of
        \Cref{theorem:SymmetricPowerOfSheafIsTamelyRamified} the field is assumed algebraically closed before the present
        theorem is invoked.
    \end{enumerate}
\end{rem*}

We now formulate the formal--local invariance precisely. Fix $P \in C(k)$ lying in the support of $\mdls$, set $\ndls = dP$,
and fix a uniformizer $x$ of $\Oo_{C,P}$; it induces an isomorphism $\widehat{\Oo_{C,P}} \cong k[[x]]$.
The scheme $\Spec(k[[x]])$ has two points --- the closed point, mapping to $P$, and the generic point, mapping to the generic
point of $C$ --- so the punctured formal disc $\Spec(k((x)))$ maps to $U$, and we denote by
$$\cP_x := \cP \times_U \Spec(k((x)))$$
the corresponding restriction of $\cP$; it is a $G$-torsor over $\Spec(k((x)))$.
Let $e_1, \dots, e_d \in k[[x_1, \dots, x_d]]$ be the elementary symmetric polynomials in $x_1, \dots, x_d$, and set
$$ V = \Spec\left(k[[e_1, \dots, e_d]][e_d^{-1}]\right), \qquad W = \Spec\left(k[[x_1, \dots, x_d]][(x_1 \cdots x_d)^{-1}]\right).$$
The inclusion $k[[e_1, \dots, e_d]] \subset k[[x_1, \dots, x_d]]$ induces a finite flat morphism $\hat{r}: W \to V$
(note that inverting $e_d = x_1 \cdots x_d$ inverts each $x_i$), and for each $i$ we let
$$ q_i : W \to \Spec(k((x))), \qquad x \mapsto x_i.$$
The group $S_d$ acts on $W$ over $V$ by permuting $x_1, \dots, x_d$, and $q_{\sigma(i)} = q_i \circ \sigma$.

\begin{lemma}[Formal--local invariance]\label{lemma:formal_local_invariance:reorg}
    In the above notation:
    \begin{enumerate}
        \item The choice of $x$ induces an isomorphism $\widehat{\Oo_{C^{(d)}, Z_\ndls}} \cong k[[e_1, \dots, e_d]]$ carrying
        the maximal ideal to $(e_1, \dots, e_d)$, such that the induced morphism $c: \Spec(k[[e_1, \dots, e_d]]) \to C^{(d)}$
        is flat and satisfies $c^{-1}(U^{(d)}) = V$; we write $c_V : V \to U^{(d)}$ for the restriction.
        \item There is an induced isomorphism $\widehat{\Oo_{X_\ndls, \eta_\ndls}} \cong \kappa[[e_d]]$, where
        $\kappa = k(\frac{e_1}{e_d}, \dots, \frac{e_{d-1}}{e_d})$. In particular
        $\widehat{K}_{\eta_\ndls} := \Frac(\widehat{\Oo_{X_\ndls, \eta_\ndls}}) \cong \kappa((e_d))$, and the canonical
        morphism $\Spec(\widehat{K}_{\eta_\ndls}) \to U^{(d)}$ factors as
        $\Spec(\kappa((e_d))) \to V \xrightarrow{c_V} U^{(d)}$, where the first map does not depend on $(C, \cP)$.
        \item There is a canonical isomorphism of $G$-torsors over $V$
        $$ \cP^{[d]} \times_{U^{(d)}} V \;\cong\; \Xi \backslash \left( q_1^{-1}\cP_x \times_W \cdots \times_W q_d^{-1}\cP_x \right).$$
    \end{enumerate}
    Consequently, the $G$-torsor $\cP^{[d]}|_{\Spec(\widehat{K}_{\eta_\ndls})}$ --- and hence whether the ramification of
    $\cP^{[d]}$ at $\eta_\ndls$ is bounded by $r$ (resp.\ tame, resp.\ unramified) --- depends only on $d$ and on the
    isomorphism class of the $G$-torsor $\cP_x$ over $k((x))$, and not otherwise on the pair $(C, \cP)$.
\end{lemma}

\begin{proof}
    \textbf{(1)} Since $P$ is a $k$-rational point of the smooth curve $C$, the uniformizer $x$ induces
    $\widehat{\Oo_{C,P}} \cong k[[x]]$. As completion at rational points takes products of $k$-varieties to completed tensor
    products, we get $\widehat{\Oo_{C^d, (P, \dots, P)}} \cong k[[x]] \widehat{\otimes}_k \cdots \widehat{\otimes}_k k[[x]] = k[[x_1, \dots, x_d]]$.
    Write $A = \Oo_{C^{(d)}, Z_\ndls}$ and $B = \Oo_{C^d, (P, \dots, P)}$.
    Since $k$ is algebraically closed, a closed point $(Q_1, \dots, Q_d) \in C^d$ maps to $Z_\ndls = dP$ under the finite
    projection $r: C^d \to C^{(d)}$ if and only if $Q_1 = \dots = Q_d = P$; hence $(P, \dots, P)$ is the unique point of
    $C^d$ over $Z_\ndls$, $B = (r_* \Oo_{C^d})_{Z_\ndls}$ is a finite $A$-module, and $A = B^{S_d}$ (formation of invariants
    commutes with localization at invariant primes).
    Since $\mathfrak{m}_B$ is the unique maximal ideal of $B$ over $\mathfrak{m}_A$, we have
    $\mathfrak{m}_B^N \subset \mathfrak{m}_A B \subset \mathfrak{m}_B$ for some $N$, so the $\mathfrak{m}_A$-adic and
    $\mathfrak{m}_B$-adic topologies on $B$ agree, and therefore $B \otimes_A \hat{A} \cong \hat{B} \cong k[[x_1, \dots, x_d]]$.
    Now $A = B^{S_d}$ is the kernel of $B \to \prod_{\sigma \in S_d} B$, $b \mapsto (\sigma(b) - b)_\sigma$; as $\hat{A}$ is
    flat over $A$, tensoring with $\hat{A}$ preserves this kernel, so
    $$\hat{A} \cong (B \otimes_A \hat{A})^{S_d} \cong k[[x_1, \dots, x_d]]^{S_d}.$$
    Finally, $k[[x_1, \dots, x_d]]^{S_d} = k[[e_1, \dots, e_d]]$: the ring $k[x_1, \dots, x_d]$ is a finite free module over
    $k[e_1, \dots, e_d]$, the same topological argument identifies $k[[x_1, \dots, x_d]]$ with
    $k[x_1, \dots, x_d] \otimes_{k[e_1, \dots, e_d]} k[[e_1, \dots, e_d]]$, and taking $S_d$-invariants --- again a kernel,
    preserved by the flat base change $k[e_1, \dots, e_d] \to k[[e_1, \dots, e_d]]$ --- yields $k[[e_1, \dots, e_d]]$.
    Under these identifications the maximal ideal of $\hat{A}$ is $(e_1, \dots, e_d)$, and $c$ is flat as the composition of
    the flat morphisms $\Spec(\hat{A}) \to \Spec(A) \to C^{(d)}$.

    It remains to check that $c^{-1}(C^{(d)} \setminus U^{(d)}) = V(e_d)$ as sets. We have
    $C^{(d)} \setminus U^{(d)} = \bigcup_{Q \in \mathrm{Supp}(\mdls)} Z_Q$ with $Z_Q = \{D : Q \in \mathrm{Supp}(D)\}$, as in
    the proof of \Cref{lemma:torsor_unramified_codim_one:reorg}. For $Q \neq P$ the divisor $Z_Q$ does not contain the point
    $Z_\ndls = dP$, so its ideal is not contained in $\mathfrak{m}_A$, and $c^{-1}(Z_Q) = \emptyset$. For $Q = P$, pull back
    along the finite surjective morphism $\Spec(\hat{B}) \to \Spec(\hat{A})$: the preimage of $Z_P$ in $C^d$ is
    $\bigcup_i p_i^{-1}(P)$, whose trace on $\Spec(\hat{B}) = \Spec(k[[x_1, \dots, x_d]])$ is
    $V(x_1) \cup \dots \cup V(x_d) = V(x_1 \cdots x_d) = V(e_d)$; by surjectivity, $c^{-1}(Z_P) = V(e_d)$ as claimed. Hence
    $c^{-1}(U^{(d)}) = V$, and likewise $\Spec(\hat{B}) \times_{C^d} U^d = W$.

    \textbf{(2)} Blowing up commutes with flat base change, and the center $Z_\ndls$ pulls back under $c$ to the closed
    point $V(e_1, \dots, e_d)$; hence
    $$\hat{Y} := \Bl_{(e_1, \dots, e_d)}\left(\Spec k[[e_1, \dots, e_d]]\right) \cong X_\ndls \times_{C^{(d)}} \Spec(\hat{A}),$$
    and the exceptional divisor $\hat{E} \subset \hat{Y}$ is the base change of $E_\ndls$ along
    $\Spec(\hat{A}/\hat{\mathfrak{m}}) = \Spec(\kappa(Z_\ndls))$, i.e.\ the projection $\mathrm{pr}: \hat{Y} \to X_\ndls$
    restricts to an isomorphism $\hat{E} \cong E_\ndls \cong \bP_k^{d-1}$.
    Let $\hat{\eta}$ be the generic point of $\hat{E}$. Then $\mathrm{pr}(\hat\eta) = \eta_\ndls$, and the induced local
    homomorphism of discrete valuation rings $\Oo_{X_\ndls, \eta_\ndls} \to \Oo_{\hat{Y}, \hat{\eta}}$ carries a uniformizer
    to a uniformizer (the ideal of $E_\ndls$ pulls back to the ideal of $\hat{E}$) and induces an isomorphism of residue
    fields (both are the function field of $\bP_k^{d-1}$); hence it induces an isomorphism on completions.
    The completion of $\Oo_{\hat{Y}, \hat{\eta}}$ is computed exactly as for the blowup of affine space at the origin, which
    we carry out in detail below: on the chart $u_d \neq 0$ we have $e_i = \frac{u_i}{u_d} e_d$, the local ring at
    $\hat\eta$ is $\left(\hat{A}[\frac{u_1}{u_d}, \dots, \frac{u_{d-1}}{u_d}]\right)_{(e_d)}$, its residue field is
    $\kappa = k(\frac{u_1}{u_d}, \dots, \frac{u_{d-1}}{u_d}) = k(\frac{e_1}{e_d}, \dots, \frac{e_{d-1}}{e_d})$, and its
    completion is $\kappa[[e_d]]$.
    Finally, $e_d$ is invertible in $\widehat{K}_{\eta_\ndls} = \kappa((e_d))$, so the composite
    $\Spec(\widehat{K}_{\eta_\ndls}) \to \Spec(\Oo_{\hat{Y}, \hat\eta}) \to \hat{Y} \to \Spec(\hat{A})$ factors through $V$,
    giving the asserted factorization of $\Spec(\widehat{K}_{\eta_\ndls}) \to U^{(d)}$ through $c_V$.

    \textbf{(3)} Let $h : \cP^{\times_k d} \to U^{(d)}$ denote the canonical morphism; it is finite, being the composition
    of the finite étale morphism $\cP^{\times_k d} \to U^d$ with the finite morphism $U^d \to U^{(d)}$. By
    \Cref{cor:sym-power-quotient},
    $\cP^{[d]} = \Xi \backslash \cP^{\times_k d} = \underline{\Spec}_{U^{(d)}}\left((h_* \Oo_{\cP^{\times_k d}})^{\Xi}\right)$.
    Pushforward along the affine morphism $h$ commutes with the flat base change $c_V : V \to U^{(d)}$, and the formation of
    $\Xi$-invariants is a kernel, hence also commutes with flat base change; therefore
    $$ \cP^{[d]} \times_{U^{(d)}} V \cong \Xi \backslash \left(\cP^{\times_k d} \times_{U^{(d)}} V\right).$$
    By step (1), $U^d \times_{U^{(d)}} V \cong W$, so
    $\cP^{\times_k d} \times_{U^{(d)}} V \cong \cP^{\times_k d} \times_{U^d} W$.
    For each $i$, the composition $W \to \Spec(k[[x_1, \dots, x_d]]) \to \Spec(k[[x]]) \to C$ --- the middle map being the
    completion of the projection $p_i : C^d \to C$, i.e.\ $x \mapsto x_i$ --- factors through $\Spec(k((x))) \to U$, since
    $x_i$ is invertible on $W$; this composite is precisely $q_i$, and it commutes with $p_i : U^d \to U$. Since
    $\cP^{\times_k d} = p_1^{-1}\cP \times_{U^d} \cdots \times_{U^d} p_d^{-1}\cP$, base changing to $W$ gives an isomorphism
    $$ \cP^{\times_k d} \times_{U^d} W \cong q_1^{-1}\cP_x \times_W \cdots \times_W q_d^{-1}\cP_x, $$
    which is equivariant for the actions of $G^d$ (in particular of $\KG$) and of $S_d$ --- the latter acting on $W$ by
    permuting the $x_i$, compatibly with $q_{\sigma(i)} = q_i \circ \sigma$. Passing to $\Xi$-quotients yields the displayed
    isomorphism.

    For the final statement: the schemes $V$ and $W$, the maps $q_i$ and $\hat{r}$, the $\Xi$-action, and the point
    $\Spec(\kappa((e_d))) \to V$ of (2) are all constructed from $d$ alone; by (3), the torsor
    $\cP^{[d]}|_{\Spec(\widehat{K}_{\eta_\ndls})}$ is the pullback along $\Spec(\kappa((e_d))) \to V$ of a torsor constructed
    from $\cP_x$ and these data. Since the ramification at $\eta_\ndls$ is by definition read off from
    $\cP^{[d]}|_{\Spec(\widehat{K}_{\eta_\ndls})}$ together with the discrete valuation ring $\kappa[[e_d]]$
    (\Cref{definition:tamely_ramified_in_codim_1:reorg}, \Cref{definition:local_ramification_boundness:reorg}), it depends only on
    $(\cP_x, d)$.
\end{proof}

\begin{cor}[Reduction to $\bP^1$]\label{cor:reduction_to_P1:reorg}
    Let $C, C'$ be smooth projective geometrically connected curves over the algebraically closed field $k$, let
    $\cP \to U = C \setminus \mdls$ and $\cP' \to U' = C' \setminus \mdls'$ be $G$-torsors, let
    $P \in \mathrm{Supp}(\mdls)(k)$ and $P' \in \mathrm{Supp}(\mdls')(k)$, and let $x, x'$ be uniformizers at $P, P'$
    respectively. Assume that $\cP_x \cong \cP'_{x'}$ as $G$-torsors, compatibly with the isomorphism
    $k((x)) \cong k((x'))$, $x \mapsto x'$.
    Then for every $d \geq 1$ and every $r$, the ramification of $\cP^{[d]}$ at $\eta_{dP}$ is bounded by $r$ if and only if
    the ramification of $\cP'^{[d]}$ at $\eta_{dP'}$ is; in particular, the one is unramified (resp.\ tamely ramified) at
    $\eta_{dP}$ if and only if the other is at $\eta_{dP'}$.
\end{cor}
\begin{proof}
    Immediate from \Cref{lemma:formal_local_invariance:reorg}: both $\cP^{[d]}|_{\Spec(\widehat{K}_{\eta_{dP}})}$ and
    $\cP'^{[d]}|_{\Spec(\widehat{K}_{\eta_{dP'}})}$ are identified, compatibly with the discrete valuation ring
    $\kappa[[e_d]]$, with the pullback along $\Spec(\kappa((e_d))) \to V$ of
    $\Xi \backslash (q_1^{-1}\cP_x \times_W \cdots \times_W q_d^{-1}\cP_x)$.
\end{proof}

\begin{lemma}[Realization over $\bG_m$]\label{lemma:realization_over_Gm:reorg}
    Let $k$ be algebraically closed of characteristic $p > 0$, and let $\cQ$ be a $G$-torsor over $\Spec(k((x)))$.
    \begin{enumerate}
        \item If $\cQ$ is tamely ramified, then there exist an integer $e$ prime to $p$, an injective homomorphism
        $\iota : \Z/e\Z \hookrightarrow G$, and an integer $a$ prime to $e$, such that
        $$\cP' := \iota_*\left( \Spec\left(k[x, x^{-1}][T]/(T^e - x^a)\right) \right)$$
        is a $G$-torsor over $\bG_m$, tamely ramified at $0$ and $\infty$, with $\cP'_x \cong \cQ$.
        Here $T^e = x^a$ is regarded as a $\Z/e\Z$-torsor via a fixed primitive $e$-th root of unity $\zeta_e \in k$, and
        $\iota_* \cQ' := (\cQ' \times G)/(\Z/e\Z)$ denotes the induced $G$-torsor.
        \item If $G = \Z/p\Z$ and $\cQ$ has ramification bounded by $d$ (\Cref{definition:local_ramification_boundness:reorg}),
        then there exists $f \in x^{-1}k[x^{-1}]$, either $f = 0$ or
        $f = cx^{-m} + a_{-m+1}x^{-m+1} + \dots + a_{-1}x^{-1}$ with $c \neq 0$, $p \nmid m$ and $m < d$, such that
        $$\cP' := \Spec\left(k[x, x^{-1}][X]/(X^p - X - f)\right)$$
        is a $\Z/p\Z$-torsor over $\bG_m$ with $\cP'_x \cong \cQ$. Moreover $\cP'$ extends to a torsor over
        $\bP^1_k \setminus \{0\}$ --- in particular it is unramified at $\infty$ --- and it has ramification bounded by $d$
        at $0$.
    \end{enumerate}
    In both cases $\cP'_x$ denotes the restriction of $\cP'$ to $\Spec(k((x)))$, as before.
\end{lemma}
\begin{proof}
    \textbf{(1)} Since $k$ is algebraically closed, the residue field of $k((x))$ is algebraically closed, so $k((x))$
    admits no non-trivial unramified extensions, and its tamely ramified extensions are exactly the Kummer extensions
    $k((x^{1/e}))$ with $p \nmid e$ (\Cref{theorem:kummer_ramification:reorg}). Hence the tame quotient $G^t$ of
    $\Gal(k((x))^{sep}/k((x)))$ is procyclic, $G^t \cong \varprojlim_{p \nmid e} \mu_e(k) \cong \prod_{\ell \neq p} \Z_\ell$;
    fix a topological generator $\gamma$.
    Let $\rho: \Gal(k((x))^{sep}/k((x))) \to G$ be the continuous homomorphism classifying $\cQ$. Tameness means that
    $\rho$ kills the wild inertia subgroup, so $\rho$ factors through $G^t$ and is determined by $g := \rho(\gamma)$, whose
    order $e$ is prime to $p$. Let $\iota: \Z/e\Z \xrightarrow{\sim} \langle g \rangle \subset G$ be given by
    $1 \mapsto g$, so that $\rho = \iota \circ \chi$ for a surjective character $\chi : G^t \to \Z/e\Z$.
    On the other hand, for $a \in \Z$ the equation $T^e = x^a$ defines a $\mu_e \cong \Z/e\Z$-torsor over
    $\Spec(k((x)))$, and by Kummer theory the classes of $x^a$, $a \in \Z/e\Z$, exhaust
    $\mathrm{Hom}_{\mathrm{cont.}}(G^t, \Z/e\Z) \cong \Z/e\Z$, the class of $x^a$ being surjective precisely when
    $\gcd(a, e) = 1$. Choose $a$ with the class of $x^a$ equal to $\chi$; then $\cQ \cong \iota_*(T^e = x^a)$ over
    $k((x))$. The same equation over $k[x, x^{-1}]$ defines a $\Z/e\Z$-torsor over $\bG_m$ (étale since $e$ is invertible
    in $k$), so $\cP' := \iota_*\left(\Spec(k[x,x^{-1}][T]/(T^e - x^a))\right)$ is a $G$-torsor over $\bG_m$ with
    $\cP'_x \cong \cQ$; it is tamely ramified at $0$ and at $\infty$ since its local characters there factor through the
    tame quotients.

    \textbf{(2)} If $\cQ$ is trivial, take $f = 0$. Otherwise, since $\Z/p\Z$ is simple and $k((x))$ admits no non-trivial
    unramified extensions, $\cQ = \Spec(F)$ for a cyclic extension $F/k((x))$ of degree $p$ together with an isomorphism
    $\Gal(F/k((x))) \cong \Z/p\Z$. By Artin--Schreier theory such data correspond to non-zero classes in
    $k((x))/\wp(k((x)))$, where $\wp(u) = u^p - u$, via $g \mapsto k((x))[X]/(X^p - X - g)$. We claim that every class
    admits a representative $f$ as in the statement:
    \begin{itemize}
        \item $\wp$ is surjective on $k[[x]]$: given $h = \sum_{i \geq 0} b_i x^i$, solve $\wp(u) = h$ with
        $u = \sum_{i \geq 0} u_i x^i$ by choosing $u_0$ with $u_0^p - u_0 = b_0$ (possible as $k$ is algebraically closed)
        and setting recursively $u_i = u_{i/p}^p - b_i$, with the convention $u_{i/p} = 0$ when $p \nmid i$. Hence, we may
        choose a representative in $x^{-1}k[x^{-1}]$.
        \item If the representative has pole order $m$ divisible by $p$, with leading term $cx^{-m}$, then subtracting
        $\wp(c^{1/p}x^{-m/p}) = cx^{-m} - c^{1/p}x^{-m/p}$ strictly decreases the pole order. Iterating, we reach a
        representative $f$ which is either $0$ (impossible, as the class is non-zero) or has pole order $m$ prime to $p$.
    \end{itemize}
    By \Cref{theorem:artin_schreier_ramification:reorg}, the unique ramification jump of $F = k((x))[X]/(X^p - X - f)$ is at $m$,
    so the ramification of $\cQ$ is bounded by $d$ if and only if $m < d$.
    Finally, $f \in k[x, x^{-1}]$, so the same equation defines a finite étale $\Z/p\Z$-cover $\cP' \to \bG_m$ (a torsor via
    $X \mapsto X + 1$) with $\cP'_x \cong \cQ$; and writing $y = x^{-1}$, we have $f \in y\,k[y]$, so $X^p - X - f$ defines
    a finite étale cover of $\mathbb{A}^1_k = \Spec(k[y]) = \bP^1_k \setminus \{0\}$ extending $\cP'$, whence $\cP'$ is
    unramified at $\infty$.
\end{proof}

Combining \Cref{cor:reduction_to_P1:reorg} with \Cref{lemma:realization_over_Gm:reorg}, in order to prove
\Cref{theorem:torsors_after_blowup_is_tame:reorg} for a torsor which is tamely ramified at $P$, or for a $\Z/p\Z$-torsor with
ramification bounded by $d$ at $P$, we may --- and do --- assume from now on that
$$C = \bP^1_k, \qquad \mdls = d \cdot 0 + \infty, \qquad \ndls = d \cdot 0, \qquad U = \bG_m,$$
and that $\cP \to \bG_m$ is one of the explicit torsors of \Cref{lemma:realization_over_Gm:reorg}.


\section{The symmetric-power computation: \texorpdfstring{$\Z/e\Z$, $\Z/p\Z$, $\Z/p^r\Z$}{Z/eZ, Z/pZ, Z/p^rZ}}
Having reduced to an explicit torsor over $\bG_m$, the computation now proceeds
identically in the three cases, along the following four steps.
\begin{description}
  \item[Step A (realization).] Present the local torsor $\cP_x$ by an explicit
  equation over $\bG_m$ --- Kummer $T^e = x^a$, Artin--Schreier $X^p - X = f$, or
  Artin--Schreier--Witt $\wp \vec X \ominus \vec f$ --- with the pole data of the
  defining datum controlled by the ramification bound.
  \item[Step B (invariant ring).] On $C^{\times d}$, form the $d$-fold product
  torsor and identify the ring of $\KG$-invariants of $S^{\otimes_k d}$: it is
  generated by a single invariant $Y$, namely $\prod_i X_i$ (Kummer),
  $\sum_i X_i$ (Artin--Schreier), or the Witt sum $\bigoplus_i \vec X_i$
  (Artin--Schreier--Witt).
  \item[Step C (descent to $e_i$).] Pass to $S_d$-invariants: the resulting
  defining datum is a symmetric function of $x_1, \dots, x_d$, hence a function
  of the elementary symmetric coordinates $e_1, \dots, e_d$.
  \item[Step D (regularity at $\eta$).] Because the pole order of the datum is
  $< d$, it is regular at $\eta_\ndls$; the cover $\cP^{[d]}|_{\widehat K_\eta}$
  is therefore tame (Kummer) or unramified (Artin--Schreier and
  Artin--Schreier--Witt).
\end{description}
The Kummer and Artin--Schreier cases are carried out together in
\Cref{theorem:ramification_after_blowup_P1:reorg}; the Artin--Schreier--Witt
case, which occupies the rest of the section, follows the same four steps in
Witt coordinates. \Cref{table:cyclic_computation_dictionary:reorg} collects the
three parallel calculations.

\begin{table}[h]
\centering\footnotesize
\renewcommand{\arraystretch}{1.4}
\begin{tabular}{@{}l>{\raggedright\arraybackslash}p{3.2cm}>{\raggedright\arraybackslash}p{3.2cm}>{\raggedright\arraybackslash}p{3.5cm}@{}}
\hline
 & \textbf{Kummer} & \textbf{Artin--Schreier} & \textbf{Artin--Schreier--Witt} \\
\hline
A: equation over $\bG_m$ & $T^e = x^a$ & $X^p - X = f$ & $\wp \vec X \ominus \vec f$ \\
B: invariant $Y$ & $\prod_i X_i$ & $\sum_i X_i$ & $\bigoplus_i \vec X_i$ \\
B: $(S^{\otimes d})^{\KG}$ & $R^{\otimes d}[Y]/(Y^e - \prod_i f_i)$ & $R^{\otimes d}[Y]/(Y^p - Y - \sum_i f(x_i))$ & $R^{\otimes d}[\vec Y]/(\wp \vec Y \ominus \vec\alpha)$ \\
C: datum in $e_i$ & $\prod_i f_i = \tilde F(e)$ & $\sum_i f(x_i) \in k(\tfrac{e_i}{e_d})$ & $\vec\alpha \in W_r\!\big(k(\tfrac{e_i}{e_d})\big)$ \\
D: bound at $\eta$ & pole $< d$ & pole $m < d$ & $p^{\,r-1-j} m_j < d$ \\
D: conclusion & tame & unramified & unramified \\
\hline
\end{tabular}
\caption{The three symmetric-power computations in parallel (Steps A--D).}
\label{table:cyclic_computation_dictionary:reorg}
\end{table}
We start by analyzing the blowup $\tilde{X} = \text{Bl}_0(X)$ of $X = \mathbb{A}_k^n$
where $0 \in X$ be the origin. 
Recall that $\tilde{X} \subset X \times_k \mathbb{P}^{n-1}_k$ is defined by the equations $x_i u_j = x_j u_i$, where $[u_1 : \dots : u_n]$ are the homogeneous coordinates of 
$\mathbb{P}^{n-1}_k$.
The exceptional divisor $E \subset \tilde{X}$ is the fiber over the origin, $E = \{ (0, [u_1 : \dots : u_n]) \}$, which is of codimension 1 in $\tilde{X}$. 
Let $\eta \in E$ be the generic point of $E$, and let $R = \mathcal{O}_{\tilde{X}, \eta}$ be the associated local ring. 
This ring $R$ is a discrete valuation ring (DVR) with fraction field $K = K(X) = k(x_1, \dots, x_n)$.

On the affine chart $U_1$ where $u_1 \neq 0$, we have $x_i = \frac{u_i}{u_1} x_1$. 
The coordinate ring is:$$\mathcal{O}_{\tilde{X}}(U_1) = k\left[x_1, \frac{u_2}{u_1}, \dots, \frac{u_n}{u_1}\right]$$
In this chart, the generic point $\eta$ corresponds to the prime ideal $\mathfrak{p}_1 = (x_1)$. Thus, the local ring is $R = k[x_1, \frac{u_2}{u_1}, \dots, \frac{u_n}{u_1}]_{(x_1)}$. 
The residue field is $\kappa(\eta) = k(\frac{u_2}{u_1}, \dots, \frac{u_n}{u_1})$, and the completion of $R$ with respect to its maximal ideal is:
$$\hat{R} = \kappa(\eta)[[x_1]]$$
In this local ring, $x_1$ is a uniformizer. Note that any $x_i$ (for $i > 1$) can also serve as a uniformizer, 
as $x_i = (\frac{u_i}{u_1}) x_1$ and $\frac{u_i}{u_1}$ is a unit in $R$.

For any monomial $M = x_1^{a_1} \dots x_n^{a_n} \in K$, we can write:
$$M = x_1^{\sum a_i} \left( \frac{u_2}{u_1} \right)^{a_2} \dots \left( \frac{u_n}{u_1} \right)^{a_n}$$
Since the term in parentheses is a unit in $R$, the valuation $\nu_E$ associated with $E$ satisfies:
$$\nu_E(M) = \sum a_i = \deg M$$
Consequently, for any polynomial $f = f_d + f_{d+1} + \dots + f_l$, where $f_i$ is the homogeneous part of degree $i$, we have $\nu_E(f) = d$ (the order of vanishing at the origin).



Our first result is:
\begin{theorem}\label{theorem:ramification_after_blowup_P1:reorg}
     Let $\cP \to \bG_m \subset \bP_k^1$ be a $G$ torsor which is either
     \begin{enumerate}
          \item tamely ramified at $0$ 
          \item totally ramified at $0$ with $G=\Z/p\Z$ and ramificaiton bounded by $d$.
     \end{enumerate}
     Then the ramification of the $G$-torsor $\cP^{[d]} \to \bG_m^{(d)}$ at $\eta_\mdls$ the generic point of $E_\mdls \subset X_{\mdls}$ is  
     \begin{enumerate}
          \item tamely ramified if $\cP$ was tamely ramified
          \item unramified if $\cP$ was totally ramified with $G=\Z/p\Z$ and ramification bounded by $d$
     \end{enumerate}
\end{theorem}
\begin{proof}
    The local ring at the generic point of the excptional divisor of the blowup of the affine space at 0 point is
     $R = k[x_d, \frac{u_1}{u_d}, \dots, \frac{u_{d-1}}{u_d}]_{(x_d)}$. 
     Where for every  $i < d$, we have $x_i = \frac{u_i}{u_d} x_d$ are all uniformizers.
     The residue field is $\kappa(\eta) = k(\frac{u_1}{u_d}, \dots, \frac{u_{d-1}}{u_d})$
     and the completion of $R$ with respect to its maximal ideal is:
     $$\hat{R} = \kappa(\eta)[[x_d]]$$
     In our situation, when we take symmetric product of the affine space, the situation is similiar with different coordinates
     if we let $e_1, \dots, e_d$ be the symmetric polynomials in $x_1, \dots x_d$ then:
     The local ring is 
     $R = k[e_d, \frac{u_2}{u_d}, \dots, \frac{u_{d-1}}{u_d}]_{(e_d)}$. 
     $e_i = \frac{u_i}{u_d} e_d$ are all uniformizers.
     The residue field being:
     $\kappa(\eta) = k(\frac{u_1}{u_d}, \dots, \frac{u_{d-1}}{u_d})$
     and the completion: $\hat{R} = \kappa(\eta)[[s_d]]$
     Note that from  $e_i = \frac{u_i}{u_d} e_d$
     We get  $\frac{e_i}{e_d} = \frac{u_i}{u_d}$ in the fracion field. hence

     \begin{equation}\label{eq:complete_field_at_generic_point:reorg}
          \hat{K}= k(\frac{e_1}{e_d}, \dots, \frac{e_{d-1}}{u_d})((e_d))
     \end{equation}

     We compute directly the extesntion of complete valued fields over the complete valued field at the generic point. 
     Note that by \Cref{cor:reduction_to_P1:reorg} and \Cref{lemma:realization_over_Gm:reorg}
     we can assume $\cP = \Spec k[x, x^{-1}][X]/(X^n - a)$ for $a \in k[x, x^{-1}]$ or $\cP = \Spec k[x, x^{-1}][X]/(X^p + X - f(x, x^{-1}))$
     where $f(x,x^{-1}) = c x^{-m} + a_{-m + 1} x^{-m + 1} + \cdots a_{-1}x^{-1} + a_0 = c x^{-m} + f_{-m+1}(x^{-1})$ where $m < d$.
     
     We deal with each case separately. 
     
     \textbf{Artin-Schreier Extensions:}

     Set $R=k[x, x^{-1}]$, and $S = \Spec k[x, x^{-1}][X]/(X^p + X - f(x^{-1}))$
     we have 
     $$
          \begin{aligned}
          R^{\otimes_k d} &= k[x_1, x_1^{-1}, \dots, x_d, x_d^{-1}] \\
          S^{\otimes_k d} &= k[x_1, x_1^{-1}, \dots, x_d, x_d^{-1}][X_1, \dots, X_d] / (X_1^p - X_1 - f(x_1), \dots, X_d^p - X_d - f(x_d)) \\
          &= R^{\otimes_k d}[X_1, \dots, X_d] / (X_1^p - X_1 - f(x_1), \dots, X_d^p - X_d - f(x_d))
          \end{aligned}
     $$

     Next, we want to understand the ring corresponding to $p_1^{-1}(\cP) \otimes \dots \otimes p_d^{-1}(\cP)$ on $C^{d}$ - 
     the $d$'th-product of the $G=\Z/p\Z$-torsors  $p_1^{-1}(\cP), \dots, p_d^{-1}(\cP)$ on $U^{d}$. 
     It corresponds to qoutient:
     $\left(p_1^{-1}(\cP) \times \dots \times p_d^{-1}(\cP)\right) / G^{d-1}$ where the action of $G^{d-1}$ on the product is:
     $$(g_1, \dots, g_{d-1}) \cdot (p_1, p_2, \dots, p_{d-1}, p_{d}) = (g_1(p_1), g_1^{-1}g_2(p_2), \dots, 
     g_{d-2}^{-1}g_{d-1}(p_{d-1}), g_{d-1}^{-1}(p_d))$$
     
     The affine ring corrsponding to the torsor product is $\left(S^{\otimes_k d}\right)^{G^{d-1}}$
    
     Recall that the action of $g \in G=\Z/p\Z$ on $X$ is $g(X) = X + g$ ($g$ correspond to a number $0 \leq g \leq p-1)$.
     So, the action of $(g_1, \dots, g_{d-1})$ on the generators $(X_1, X_2, \dots, X_{d-1}, X_{d})$
     Is $X_1 \mapsto X_1 + g_1$, $X_i \mapsto X_i - g_{i-1} + g_i$ for $1<i<d$ and $X_d \mapsto X_d - g_{d-1}$.
     So we see that $Y=X_1 + \dots + X_d$ is invariant.
     Moreover $Y^P - Y - \sum_{i=1}^d f(x_i)=0$ is irreducible degree $p$ equation for $Y$,
     Since we are quotienting a rank $p^{d}$ extesntion by a group of order $p^{d-1}$ the resulting invaraint subring must
     have rank $p$ over $R^{\otimes_k d}$, So we conclude:
     \[
          \left(S^{\otimes_k d}\right)^{G^{d-1}} \cong R^{\otimes_k d}[Y] / (Y^p - Y - \sum_{i=1}^d f(x_i))
     \]  
     
     The group $S_d$ acts on 
     $ R^{\otimes_k d} = k[x_1^{\pm 1}, x_2^{\pm 1}, \dots, x_d^{\pm 1}]$
      by permuting the variables $\{x_i\}_{i=1}^d$.
     And since $Y=\sum_i^d X_i$ it leaves $Y$ invaraint. 
     The invaraint subring $(R^{\otimes_k d})^{S_d}$ is simply $k[e_1, e_2, \dots, e_d, e_d^{-1}]$
     where $\{ e_i \}$  are the symmetric polynomials in $x_1, \dots, x_d$:
     $$e_k = \sum_{1 \le j_1 < j_2 < \dots < j_k \le d} x_{j_1} x_{j_2} \dots x_{j_k}$$
     
     To find $\left(R^{\otimes_k d}[Y] / (Y^p - Y - \sum_{i=1}^d f(x_i))\right)^{S_d}$ its enough to express
     $\sum_{i=1}^d f(x_i)$ in $e_1, \dots, e_d$, this can be done with the newton polynomials, moreover, we claim the following:

%% [REVIEW: unlabeled -- used only in the immediately following 'Finally, restricting...' step; fine inline, but a label would make the Kummer case's parallel argument citable]
     \begin{lemma}
          Let $f(x) = c x^{-m} + a_{-m + 1} x^{-m + 1} + \cdots a_{-1}x^{-1} + a_0$, 
          define  $\alpha(x_1, ..., x_d) = \sum_1^d f(x_i)$, and deonte by $e_1, ..., e_d$ the elementary symmetric polynomials 
          in $x_1, ..., x_d$. If $m < d$ then $\alpha(x_1, ..., x_d) \in k(e_1/e_d, e_2/e_d, ..., e_{d-1}/e_d)$
     \end{lemma}
     \begin{proof}
          Changing variables $y_i=x_i^{-1}$ for each $i \in \{1, \dots, d\}$
          We get 
          $$\alpha = \sum_{i=1}^d f(x_i) = \sum_{i=1}^d \left( c y_i^m + a_{-m+1} y_i^{m-1} + \dots + a_{-1} y_i + a_0 \right)$$
          Rearranging the sums, we get
          $$\alpha = c \sum_{i=1}^d y_i^m + a_{-m+1} \sum_{i=1}^d y_i^{m-1} + \dots + a_{-1} \sum_{i=1}^d y_i + d a_0$$
          Let $p_k(y_1, \dots, y_d) = \sum_{i=1}^d y_i^k$ be the $k$-th power sum symmetric polynomial. 
          The expression for $\alpha$ is a linear combination of these power sums:
          $$\alpha = c p_m(y) + a_{-m+1} p_{m-1}(y) + \dots + a_{-1} p_1(y) + d a_0$$
         
          According to the \textit{Fundamental Theorem of Symmetric Polynomials}, any symmetric polynomial in $y_1, \dots, y_d$ 
          can be expressed as a polynomial in the elementary symmetric polynomials $e_k(y_1, \dots, y_d)$. 
          Since $m < d$, $\alpha$ is a polynomial in $e_1(y), e_2(y), \dots, e_m(y)$. ($y=(y_1, \dots, y_d)$)
          The elementary symmetric polynomials in $y_i = 1/x_i$ are related to the elementary symmetric polynomials in $x_i$ as follows: 
          $$e_k(y_1, \dots, y_d) = \sum_{1 \le i_1 < \dots < i_k \le d} \frac{1}{x_{i_1} \dots x_{i_k}} = \frac{\sum_{1 \le j_1 < \dots < j_{d-k} \le d} x_{j_1} \dots x_{j_{d-k}}}{x_1 x_2 \dots x_d}$$
          Thus,
          $$e_k(y_1, \dots, y_d) = \frac{e_{d-k}(x_1, \dots, x_d)}{e_d(x_1, \dots, x_d)}$$
          which concludes the proof.
     \end{proof}

Finally, restricting $\cP^{[d]}$ to $\spec {\hat{K}}$ we get by \Cref{eq:complete_field_at_generic_point:reorg}
and \Cref{theorem:artin_schreier_ramification:reorg} the result. (that $\cP$ is unramified at the generic point of the exceptional divisor of the blowup).

\textbf{Kummer Extensions:}
Few things are different in that case,

     Set $R=k[x, x^{-1}]$, and $S=R[X]/(X^n - f)$ where $f= f(x, 1/x) \in R$
     In this case we have $char k = p$ and $gcd(p, n)=1$.

     We have 
     $$
          \begin{aligned}
          R^{\otimes_k d} &= k[x_1, x_1^{-1}, \dots, x_d, x_d^{-1}] \\
          S^{\otimes_k d} &= k[x_1, x_1^{-1}, \dots, x_d, x_d^{-1}][X_1, \dots, X_d] / (X_1^n - f_1, \dots, X_d^n - f_d) \\
          &= R^{\otimes_k d}[X_1, \dots, X_d] / (X_1^n - f_1, \dots, X_d^n - f_d)
          \end{aligned}
     $$
     Where $f_i = f(x_i, x_i^{-1})$

     Next, we want to figure out 
      $\left(S^{\otimes_k d}\right)^{G^{d-1}}$
    

     Recall that the action of $g \in G=\Z/n\Z$ on $X$ is $g(X) = \zeta^{g} X$ ($g$ correspond to a number $0 \leq g \leq n-1)$.
     So, the action of $(g_1, \dots, g_{d-1})$ on the generators $(X_1, X_2, \dots, X_{d-1}, X_{d})$
     Is $X_1 \mapsto \zeta^{g_1} X_1$, $X_i \mapsto \zeta^{g_i - g_{i-1}}X_i$ for $1<i<d$ and $X_d \mapsto \zeta^{-g_{d-1}} X_d$.

     So we see that $Y=X_1 X_2 \dots X_d$ is invariant.
     And $Y^n - \prod_{i=1}^d f_i$ is irreducible degree $n$ equation for $Y$,
     So, like before, we conclude:
     \[
          \left(S^{\otimes_k d}\right)^{G^{d-1}} \cong R^{\otimes_k d}[Y] / (Y^n - \prod_{i=1}^d f_i)
     \]  
     
     The group $S_d$ acts on $R^{\otimes_k d}[Y] / (Y^n - \prod_{i=1}^d f_i)$ by permuting the indices.
     on the variables $x_i$, 
     On $Y=\prod_1^d X_i$ it is invaraint. 
     The invaraint subring $(R^{\otimes_k d})^{S_d}$ is simply $k[e_1, e_2, \dots, e_d, e_d^{-1}]$ like before.
     

     The polynomial $F=\prod_{i=1}^d f_i$ is symmetric in $\{x_i\}_1^d$ so it can be expressed as a polynomial
     $\tilde{F}(e_1, \dots, e_d)$ in the elementary symmetric variables. 
     Hence the qoutient ring is:
     $$\left( \frac{k[x_1^{\pm 1}, \dots, t_d^{\pm 1}][Y]}{(Y^n - \prod_{i=1}^d f_i)} \right)^{S_d} \cong 
     \frac{k[e_1, \dots, e_d, e_d^{-1}][Y]}{(Y^n - \tilde{F}(e_1, \dots, e_d))}$$
     So we see again, that  restricting $\cP^{[d]}$ to $\spec {\hat{K}}$ we get by \Cref{eq:complete_field_at_generic_point:reorg} 
     and \Cref{theorem:kummer_ramification:reorg}, that $\cP$ is tamely ramified at the generic point of the exceptional divisor of the blowup.
\end{proof}

So, we conclude 
\begin{cor}\label{cor:ramification_after_blowup_is_tame:reorg}
    In the general settings of a curve $C$, and $\cP \to U=C \setminus \mdls$ a $G$-torsor with ramification bounded by $\ndls = n P \subset \mdls$ at $P$, we have:
    \begin{enumerate}
        \item If $G=\Z/p\Z$ then $\cP^{[\deg \ndls]}$ is unramified at $\eta_\ndls$
        \item If $\cP$ is tamely ramified at $\ndls$ then $\cP^{[\deg \ndls]}$ is tamely ramified at $\eta_\ndls$
    \end{enumerate}
\end{cor}
\begin{proof}
    By the remark opening this section we may assume $k$ is algebraically closed, so that $P \in C(k)$ and
    $\deg \ndls = n$. The restriction $\cP_x$ of $\cP$ to the punctured formal disc at $P$ is tamely ramified, resp.\ a
    $\Z/p\Z$-torsor with ramification bounded by $n$, so by \Cref{lemma:realization_over_Gm:reorg} there is a torsor
    $\cP' \to \bG_m$ of the corresponding explicit form with $\cP'_x \cong \cP_x$. By \Cref{cor:reduction_to_P1:reorg}, the
    ramification of $\cP^{[\deg \ndls]}$ at $\eta_\ndls$ agrees with that of $\cP'^{[\deg \ndls]}$ at $\eta_{n \cdot 0}$,
    which was computed in \Cref{theorem:ramification_after_blowup_P1:reorg}.
\end{proof}


% =====================================================================
%  The case G = Z/p^r Z.
%  ---------------------------------------------------------------
%  Imported from the working note induction_completion/gap_closure
%  ("Closing the gaps"), adapted to the thesis conventions.  The
%  general Witt-vector background it relies on lives in the
%  Preliminaries (witt/WittVectorPrerequisites.tex); this file contains
%  only the parts specific to the theorem: the W_r-realization lemma
%  over G_m (Gap 1), the identification of the symmetric-power cover
%  (Gap 2), the degree count, and the proof itself.
% =====================================================================

Next, we generalize to the case $G = \Z/p^r\Z$. This is the third column of
\Cref{table:cyclic_computation_dictionary:reorg}, and it follows the same
Steps A--D as the Kummer and Artin--Schreier cases of
\Cref{theorem:ramification_after_blowup_P1:reorg}: the paragraph
``Realization over $\bG_m$ for Witt vectors'' below is Step~A, ``the
symmetric-power cover at $\eta_\ndls$'' is Steps~B--C, and ``the degree count''
together with the final proof is Step~D --- each carried out in Witt
coordinates using the theory of
\Cref{subsection:witt_vector_prerequisites:reorg}.

\begin{theorem}\label{theorem:ramification_after_blowup_is_tame_for_p_power:reorg}
    Assume $G=\Z/p^r\Z$. Then $\cP^{[\deg (\ndls)]}$ is unramified at $\eta_{\ndls}$.
\end{theorem}

The proof occupies the remainder of this subsection. In contrast with what the phrase
``we generalize'' may suggest, the argument is \emph{not} an induction on $r$ with
\Cref{cor:ramification_after_blowup_is_tame:reorg}(1) as its base case: it is a direct
argument, treating all $r$ at once, in which each of the three steps of the
Artin--Schreier case of \Cref{theorem:ramification_after_blowup_P1:reorg} --- realization of
the local torsor by an explicit equation over $\bG_m$, identification of the
symmetric-power cover at $\eta_\ndls$, and the degree estimate --- is replaced by its
Witt-vector analogue, using the theory collected in
\Cref{subsection:witt_vector_prerequisites:reorg}. In particular the case $r=1$ is
\emph{reproved} (as the special case of Witt vectors of length one) rather than used.
For an explanation of why the seemingly natural induction on $r$ fails --- in short:
the hypothesis bounds the ramification of $\cP$ in \emph{upper} numbering, while the
break of the top degree-$p$ layer $\cP \to \cP'$ is the largest \emph{lower}-numbering
break, which can be of size about $p^{r-1}d$ --- see \Cref{section:a_wrong_path}.

\subsubsection*{Standing reduction}

As in the proof of \Cref{cor:ramification_after_blowup_is_tame:reorg}, we may assume $k$ is
algebraically closed. We may further assume that $\cP \to U$ is \emph{totally
ramified} at $P$ with inertia equal to all of $G$: writing
$R = \widehat{\Oo_{C,P}}$, $K = \Frac(R)$, and letting
$H = \rho(G^0) \subseteq G$ be the image of the inertia subgroup under the
homomorphism $\rho \colon \Gal(K^{sep}/K) \to G$ classifying $\cP|_{\Spec K}$,
decompose $\cP \to \cP/H \to U$ as in \Cref{prop:quotient_torsors}, where
$\cP \to \cP/H$ is an $H$-torsor and $\cP/H \to U$ a $G/H$-torsor. Then
$\cP/H$ is unramified at $P$, hence extends over $P$; taking any point
$Q \in \cP/H$ over $P$ yields $\widehat{\Oo_{\cP/H, Q}} \cong \widehat{\Oo_{C,P}}$,
and since $k$ is algebraically closed the torsor $\cP/H$ is locally trivial at
$Q$. This replaces $(G, \cP \to U)$ by $(H, \cP \to \cP/H)$ with
$H \cong \Z/p^s\Z$ for some $s \leq r$; as the theorem is proved for all $r$
simultaneously, we may rename and assume $H = G$.

After this reduction, the restriction $\cQ := \cP_x$ of $\cP$ to the punctured formal
disc $\Spec k((x))$ at $P$ (with $x$ a uniformizer at $P$) is a \emph{connected}
$\Z/p^r\Z$-torsor: $\cQ = \Spec L$ for a cyclic, totally ramified extension $L/K$ of
degree $p^r$, with ramification bounded by $d := \deg \ndls$
(\Cref{definition:local_ramification_boundness:reorg}); equivalently, the largest
upper-numbering break $u$ of $L/K$ satisfies $u < d$ \emph{strictly}.

\subsubsection*{Realization over $\bG_m$ for Witt vectors}

The first step generalizes \Cref{lemma:realization_over_Gm:reorg}(2) from $\Z/p\Z$ to
$\Z/p^r\Z$: every local torsor class as above is realized by an explicit
Artin--Schreier--Witt equation over $\bG_m$ whose Witt datum is a vector of
\emph{polynomials in $x^{-1}$} with controlled pole orders. Recall from
\Cref{definition:reduced_witt_vector:reorg} the notion of a reduced vector; we add:

\begin{definition}\label{definition:reduced_polynomial_vector:reorg}
    A vector $\vec g = (g_0, \dots, g_{r-1}) \in W_r(k((x)))$ is a \textbf{reduced
    polynomial vector} if it is reduced (\Cref{definition:reduced_witt_vector:reorg}) and
    moreover every $g_j \in x^{-1}k[x^{-1}]$; then $m(g_j) = \deg_{x^{-1}} g_j$, and
    $m(g_j) = 0$ if and only if $g_j = 0$.
\end{definition}

\begin{lemma}[surjectivity of $\wp$ on integral Witt vectors]\label{lemma:wp_surjective_on_integral_witt:reorg}
    $\wp \colon W_r(k[[x]]) \to W_r(k[[x]])$ is surjective.
\end{lemma}

\begin{proof}
    Induction on $r$. For $r=1$ this is the first bullet in the proof of
    \Cref{lemma:realization_over_Gm:reorg}(2): given $h = \sum_{i \geq 0} b_i x^i$, solve
    $u^p - u = h$ coefficientwise, using that $k$ is algebraically closed for the
    constant term. For $r>1$, let $\vec h \in W_r(k[[x]])$. By induction choose
    $\vec u' \in W_{r-1}(k[[x]])$ with $\wp \vec u' = t(\vec h)$, and lift it to
    $\tilde u := (\vec u', 0) \in W_r(k[[x]])$. By
    \Cref{lemma:witt_bookkeeping:reorg}(i),
    $t\bigl(\wp \tilde u \ominus \vec h\bigr) = \wp \vec u' \ominus t(\vec h) = 0$, so
    $\wp \tilde u \ominus \vec h = V^{r-1}(g)$ for some $g \in k[[x]]$. Choose
    $w \in k[[x]]$ with $\wp w = -g$ (case $r=1$ again). Then by
    \Cref{lemma:witt_bookkeeping:reorg}(ii) and additivity of $V^{r-1}$,
    \[
    \wp\bigl(\tilde u \oplus V^{r-1}(w)\bigr)
    = \wp \tilde u \oplus V^{r-1}(\wp w)
    = \bigl(\vec h \oplus V^{r-1}(g)\bigr) \oplus V^{r-1}(-g) = \vec h. \qedhere
    \]
\end{proof}

\begin{prop}[reduced polynomial representatives]\label{prop:reduced_polynomial_representatives:reorg}
    Every class in $W_r(k((x)))/\wp W_r(k((x)))$ has a reduced polynomial
    representative: for every $\vec f \in W_r(k((x)))$ there exists
    $\vec u \in W_r(k((x)))$ such that $\vec g := \vec f \ominus \wp \vec u$ has all
    components in $x^{-1}k[x^{-1}]$ with pole orders either $0$ or prime to $p$.
\end{prop}

\begin{proof}
    Induction on $r$.

    \emph{Base case $r=1$.} This is the argument of
    \Cref{lemma:realization_over_Gm:reorg}(2): first, writing $f = f^- + f^+$ with
    $f^- \in x^{-1}k[x^{-1}]$ and $f^+ \in k[[x]]$, surjectivity of $\wp$ on $k[[x]]$
    removes $f^+$; then, as long as the leading term is $cx^{-m}$ with $p \mid m$,
    $m>0$, subtracting $\wp(c^{1/p}x^{-m/p}) = cx^{-m} - c^{1/p}x^{-m/p}$ strictly
    decreases the pole order while staying inside $x^{-1}k[x^{-1}]$; this terminates
    at $0$ or at pole order prime to $p$.

    \emph{Inductive step.} Let $\vec f \in W_r(k((x)))$. By induction there is
    $\vec u' \in W_{r-1}(k((x)))$ such that
    $\vec g' := t(\vec f) \ominus \wp \vec u'$ is a reduced polynomial vector in
    $W_{r-1}$. Lift $\vec u'$ to $\tilde u := (\vec u', 0) \in W_r(k((x)))$ and set
    $\vec h := \vec f \ominus \wp \tilde u$. By \Cref{lemma:witt_bookkeeping:reorg}(i),
    $t(\vec h) = \vec g'$, i.e.\
    \[
    \vec h = (g'_0, \dots, g'_{r-2},\, h_{r-1}), \qquad h_{r-1} \in k((x)),
    \]
    with the first $r-1$ components already reduced polynomials. It remains to repair
    the last component \emph{without touching the others}; by
    \Cref{lemma:witt_bookkeeping:reorg}(iii), adjustments of the form
    $\ominus\, \wp\bigl(V^{r-1}(z)\bigr) = \ominus\, V^{r-1}(\wp z)$ change only the
    last component, by $-\wp(z)$. So:
    \begin{itemize}
        \item Write $h_{r-1} = h^- + h^+$ with $h^- \in x^{-1}k[x^{-1}]$,
        $h^+ \in k[[x]]$, choose $w \in k[[x]]$ with $\wp w = h^+$
        (\Cref{lemma:wp_surjective_on_integral_witt:reorg} with $r=1$), and replace
        $\vec h$ by $\vec h \ominus \wp(V^{r-1}w)$: the last component becomes $h^-$.
        \item While the last component has leading term $cx^{-m}$ with $p \mid m$,
        $m>0$, replace $\vec h$ by
        $\vec h \ominus \wp\bigl(V^{r-1}(c^{1/p}x^{-m/p})\bigr)$: the last component
        changes by $-\wp(c^{1/p}x^{-m/p}) = -cx^{-m} + c^{1/p}x^{-m/p}$, so its pole
        order strictly decreases and it stays in $x^{-1}k[x^{-1}]$. This terminates at
        $0$ or at pole order prime to $p$.
    \end{itemize}
    The result is a reduced polynomial vector $\wp$-equivalent to $\vec f$.
\end{proof}

\begin{rem*}
    The individual components of a Witt vector cannot be truncated to their principal
    parts one at a time --- Witt addition is not componentwise
    (\Cref{eq:witt_addition_shape:reorg}), so a componentwise operation changes the
    Artin--Schreier--Witt class. The point of
    \Cref{prop:reduced_polynomial_representatives:reorg} is that all the adjustments are
    performed by subtracting genuine elements of $\wp W_r(k((x)))$, inside the Witt
    ring.
\end{rem*}

\begin{lemma}[Realization over $\bG_m$ for $W_r$]\label{lemma:realization_over_Gm_Wr:reorg}
    Let $k$ be algebraically closed of characteristic $p>0$ and let $\cQ$ be a
    \emph{connected} $\Z/p^r\Z$-torsor over $\Spec(k((x)))$ with ramification bounded
    by $d$ (\Cref{definition:local_ramification_boundness:reorg}). Then there is a reduced
    polynomial vector $\vec f = (f_0, \dots, f_{r-1})$, $f_j \in x^{-1}k[x^{-1}]$,
    such that, with $m_j := \deg_{x^{-1}} f_j$:
    \begin{enumerate}
        \item \label{item:realization_Wr_polebound:reorg} $p^{\,r-1-j}\, m_j < d$ for all
        $0 \leq j \leq r-1$; equivalently $m_j < d\, p^{\,j-(r-1)}$;
        \item \label{item:realization_Wr_torsor:reorg}
        $\cP' := \Spec\Bigl(k[x,x^{-1}][\vec X]\big/\bigl(\wp \vec X \ominus \vec f\bigr)\Bigr)$
        is a $\Z/p^r\Z$-torsor over $\bG_m$ with $\cP'_x \cong \cQ$;
        \item \label{item:realization_Wr_infty:reorg} $\cP'$ extends to a torsor over
        $\bP^1_k \setminus \{0\}$ --- in particular it is unramified at $\infty$ ---
        and has ramification bounded by $d$ at $0$.
    \end{enumerate}
\end{lemma}

\begin{proof}
    $\cQ = \Spec L$ with $L/K$ cyclic totally ramified of degree $p^r$ and largest
    upper break $u < d$ (\Cref{theorem:asw_theory:reorg}). Let
    $[\vec f_0] \in W_r(K)/\wp$ be its Artin--Schreier--Witt class and let $\vec f$ be
    a reduced polynomial representative
    (\Cref{prop:reduced_polynomial_representatives:reorg}).

    First, $f_0 \neq 0$: otherwise
    \Cref{cor:asw_class_order_and_first_component:reorg} would bound the order of the class
    by $p^{r-1}$, contradicting $[L:K] = p^r$. Hence $m_0 \geq 1$ and $p \nmid m_0$,
    and the Schmid--Witt break formula
    (\Cref{theorem:schmid_witt_break_formula:reorg}) applies:
    $\max_j p^{r-1-j} m_j = u < d$, which is \ref{item:realization_Wr_polebound:reorg}.

    For \ref{item:realization_Wr_torsor:reorg}: by \Cref{lemma:asw_etale_torsor:reorg} (with
    $A = k[x,x^{-1}] \ni f_j$), $\cP'$ is a $\Z/p^r\Z$-torsor over $\bG_m$; its
    restriction to $\Spec k((x))$ is the Artin--Schreier--Witt torsor of the class of
    $\vec f$, i.e.\ $\cQ$ (\Cref{theorem:asw_theory:reorg}).

    For \ref{item:realization_Wr_infty:reorg}: in the coordinate $y = x^{-1}$ at $\infty$ we
    have $f_j \in y\,k[y]$, so $\vec f \in W_r(k[y])$ and
    \Cref{lemma:asw_etale_torsor:reorg} (with $A = k[y]$) extends $\cP'$ to a torsor over
    $\A^1_k = \Spec k[y] = \bP^1_k \setminus \{0\}$; a torsor over the full local ring
    at $\infty$ is unramified there. The ramification bound at $0$ holds because
    $\cP'_x \cong \cQ$.
\end{proof}

\begin{rem*}\label{rem:realization_Wr_sanity:reorg}
    Connectedness of $\cQ$ is guaranteed in the application by the standing reduction
    (inertia equal to all of $G$). Note also the following sanity check: for every
    $j$ with $p^{\,r-1-j} \geq d$, the bound \ref{item:realization_Wr_polebound:reorg}
    forces $m_j = 0$, i.e.\ $f_j = 0$ --- all sufficiently low Witt components of a
    reduced representative vanish outright. This makes visible how strong the
    hypothesis $u < d$ is at the lower levels.
\end{rem*}

\subsubsection*{The symmetric-power cover at $\eta_\ndls$}

We now generalize the invariant-ring computation in the Artin--Schreier case of
\Cref{theorem:ramification_after_blowup_P1:reorg} to Witt coordinates. Fix
$\vec f \in W_r(R)$, $R = k[x, x^{-1}]$, and let
\[
S \;=\; R[\vec X]\big/\bigl(\wp \vec X \ominus \vec f\bigr)
\]
be the coordinate ring of the torsor $\cP' \to \bG_m$ of
\Cref{lemma:realization_over_Gm_Wr:reorg}. Write
$R^{\otimes_k d} = k[x_1^{\pm 1}, \dots, x_d^{\pm 1}]$ and
\[
S^{\otimes_k d}
= R^{\otimes_k d}[\vec X_1, \dots, \vec X_d]\big/
\bigl(\wp \vec X_1 \ominus \vec f(x_1), \dots, \wp \vec X_d \ominus \vec f(x_d)\bigr),
\]
where $\vec X_i = (X_{i,0}, \dots, X_{i,r-1})$ and $\vec f(x_i)$ is the image of
$\vec f$ under $x \mapsto x_i$. Exactly as in the proof of
\Cref{theorem:ramification_after_blowup_P1:reorg}, the $d$-fold product torsor
$p_1^{-1}\cP' \otimes^G \cdots \otimes^G p_d^{-1}\cP'$ on $U^d$ is the quotient of
$\Spec S^{\otimes_k d}$ by $\KG \cong G^{d-1}$, embedded in $G^d$ as the kernel of the
sum map via $(g_1, \dots, g_{d-1}) \mapsto (g_1, g_2 - g_1, \dots, -g_{d-1})$, and
acting (in Witt coordinates, writing $\underline g$ for the image of $g \in G$ under
$\Z/p^r\Z \cong W_r(\F_p)$, \Cref{example:witt_of_Fp:reorg}) by
\[
(g_1, \dots, g_{d-1}) \colon \quad
\vec X_1 \mapsto \vec X_1 \oplus \underline g_1, \qquad
\vec X_i \mapsto \vec X_i \oplus \underline g_i \ominus \underline g_{i-1}
\ (1 < i < d), \qquad
\vec X_d \mapsto \vec X_d \ominus \underline g_{d-1}.
\]

\begin{prop}[identification of the product torsor]\label{prop:witt_cover_identification:reorg}
    Set
    \[
    \vec Y := \bigoplus_{i=1}^{d} \vec X_i \in W_r\bigl(S^{\otimes_k d}\bigr),
    \qquad
    \vec \alpha := \bigoplus_{i=1}^{d} \vec f(x_i) \in W_r\bigl(R^{\otimes_k d}\bigr),
    \]
    where $\bigoplus$ denotes the $d$-fold Witt vector sum $\oplus$ (not a direct
    sum); in particular $\vec Y$ is a single Witt vector of length $r$. Then:
    \begin{enumerate}
        \item \label{item:witt_cover_invariance:reorg} Every component $Y_n$ of $\vec Y$ is
        $G^{d-1}$-invariant, and
        \[
        \bigl(S^{\otimes_k d}\bigr)^{G^{d-1}}
        \;\cong\;
        R^{\otimes_k d}[\vec Y]\big/\bigl(\wp \vec Y \ominus \vec \alpha\bigr)
        \]
        as $G$-torsors over $R^{\otimes_k d}$ (the residual $G = G^d/G^{d-1}$ acting
        by $\vec Y \mapsto \vec Y \oplus \underline g$).
        \item \label{item:witt_cover_symmetry:reorg} Every component $\alpha_n$ of
        $\vec \alpha$ is a symmetric Laurent polynomial, i.e.\ lies in
        $\bigl(R^{\otimes_k d}\bigr)^{S_d} = k[e_1, \dots, e_d, e_d^{-1}]$, and
        \[
        \Bigl(\bigl(S^{\otimes_k d}\bigr)^{G^{d-1}}\Bigr)^{S_d}
        \;\cong\;
        k[e_1, \dots, e_d, e_d^{-1}][\vec Y]\big/\bigl(\wp \vec Y \ominus \vec \alpha\bigr).
        \]
    \end{enumerate}
    Consequently, restricting to the complete local field
    $\widehat K_\eta = \kappa(\eta)((e_d))$,
    $\kappa(\eta) = k(e_1/e_d, \dots, e_{d-1}/e_d)$, of
    \Cref{eq:complete_field_at_generic_point:reorg},
    \[
    \cP'^{[d]}\big|_{\Spec \widehat K_\eta}
    \;\cong\;
    \Spec\; \widehat K_\eta[\vec Z]\big/\bigl(\wp \vec Z \ominus \vec \alpha\bigr).
    \]
\end{prop}

\begin{proof}
    \ref{item:witt_cover_invariance:reorg} An element
    $\sigma = (g_1, \dots, g_{d-1}) \in G^{d-1}$ acts as a ring automorphism of
    $S^{\otimes_k d}$: this is its \emph{original} action, defined on the generators
    $X_{i,n}$ by the twisted formulas above (which mix components through the carry
    polynomials of \Cref{eq:witt_addition_shape:reorg}). By functoriality of $W_r$
    (\Cref{subsection:witt_vector_prerequisites:reorg}), the induced map $W_r(\sigma)$ on
    $W_r(S^{\otimes_k d})$ is componentwise and commutes with $\oplus, \ominus$; on
    the generators it reproduces the twisted action,
    $W_r(\sigma)(\vec X_i) = \vec X_i \oplus \underline g_i \ominus \underline g_{i-1}$
    (with $\underline g_0 = \underline g_d = \vec 0$). Hence, using commutativity and
    associativity of $\oplus$, the Witt sum transforms by the telescoping constant
    \[
    W_r(\sigma)(\vec Y)
    \;=\; \bigoplus_{i=1}^{d} W_r(\sigma)(\vec X_i)
    \;=\; \vec Y \oplus \underline g_1 \oplus (\underline g_2 \ominus \underline g_1)
    \oplus \cdots \oplus (\ominus\, \underline g_{d-1})
    \;=\; \vec Y.
    \]
    Since $W_r(\sigma)$ is componentwise, this equality of $r$-tuples says exactly
    $\sigma(Y_n) = Y_n$ for every $n$, i.e.\ each $Y_n \in S^{\otimes_k d}$ is
    invariant. Since $\wp = F - \mathrm{id}$ is additive
    (\Cref{definition:asw_operator:reorg} ff.),
    \[
    \wp \vec Y = \bigoplus_{i=1}^{d} \wp \vec X_i
    = \bigoplus_{i=1}^{d} \vec f(x_i) = \vec \alpha .
    \]
    Hence there is a well-defined $R^{\otimes_k d}$-algebra homomorphism
    \[
    \phi \colon\; T := R^{\otimes_k d}[\vec Z]\big/\bigl(\wp \vec Z \ominus \vec \alpha\bigr)
    \;\longrightarrow\; \bigl(S^{\otimes_k d}\bigr)^{G^{d-1}},
    \qquad \vec Z \mapsto \vec Y .
    \]
    By \Cref{lemma:asw_etale_torsor:reorg}, $\Spec T$ is a $G$-torsor over $U^d$. The
    target is the product torsor (the contracted product of the $p_i^{-1}\cP'$ along
    the sum map $G^d \to G$), a $G$-torsor over $U^d$ by
    \Cref{prop:symmetric_power_torsor} and the surrounding discussion; its residual
    $G$-action is induced by $(g, 0, \dots, 0) \in G^d$, i.e.\
    $\vec X_1 \mapsto \vec X_1 \oplus \underline g$, under which
    $\vec Y \mapsto \vec Y \oplus \underline g$. So $\Spec \phi$ is a $G$-equivariant
    morphism of $G$-torsors over $U^d$, hence an isomorphism by
    \Cref{prop:torsor_morphism_iso}.

    \ref{item:witt_cover_symmetry:reorg} $S_d$ acts on $S^{\otimes_k d}$ by permuting the
    tensor factors, i.e.\ $x_i \mapsto x_{\sigma(i)}$,
    $\vec X_i \mapsto \vec X_{\sigma(i)}$. As in
    \ref{item:witt_cover_invariance:reorg}, each $\sigma \in S_d$ is a ring automorphism,
    so $W_r(\sigma)$ is componentwise and commutes with $\oplus$. Since $\oplus$ is
    commutative and associative, permuting the summands of a Witt sum leaves it
    unchanged; thus $\vec Y$ is $S_d$-invariant (componentwise, as before) and
    $W_r(\sigma)(\vec \alpha) = \bigoplus_i \vec f(x_{\sigma(i)}) = \vec \alpha$, so
    each $\alpha_n \in R^{\otimes_k d}$ is a symmetric Laurent polynomial, hence lies
    in $k[e_1, \dots, e_d, e_d^{-1}]$ (the $S_d$-invariants of $R^{\otimes_k d}$, as
    in the proof of \Cref{theorem:ramification_after_blowup_P1:reorg}). By
    \Cref{lemma:asw_etale_torsor:reorg}, $T$ is free over $R^{\otimes_k d}$ with basis the
    monomials $\prod_n Y_n^{a_n}$, $0 \leq a_n < p$, each of which is
    $S_d$-invariant; writing an element in this basis, it is $S_d$-invariant if and
    only if all its coefficients are. Therefore
    \[
    T^{S_d} = \bigl(R^{\otimes_k d}\bigr)^{S_d}[\vec Y]\big/\bigl(\wp \vec Y \ominus \vec \alpha\bigr)
    = k[e_1, \dots, e_d, e_d^{-1}][\vec Y]\big/\bigl(\wp \vec Y \ominus \vec \alpha\bigr).
    \]

    The final statement follows exactly as in the $r=1$ case: by definition
    $\cP'^{[d]}$ is the quotient by $\Xi = \KG \rtimes S_d$ computed above
    (\Cref{cor:sym-power-quotient}), and its restriction to $\Spec \widehat K_\eta$
    along $k[e_1, \dots, e_d, e_d^{-1}] \to \widehat K_\eta$ is the
    Artin--Schreier--Witt cover of the image of $\vec \alpha$ in
    $W_r(\widehat K_\eta)$.
\end{proof}

\begin{rem*}
    The proof via \Cref{prop:torsor_morphism_iso} replaces the rank count of the
    $r=1$ argument ($p^{rd}/p^{r(d-1)} = p^r$ together with irreducibility of the
    invariant's minimal equation); the rank count alone would still require showing
    that $\vec Y$ generates the invariant ring, which is exactly what equivariance
    gives for free. Note also that a pure restriction argument on $H^1$'s (``the
    class of the product torsor is the sum of the pulled-back characters'')
    identifies the cover over $k(x_1, \dots, x_d)$ but does \emph{not} by itself
    descend it to the symmetric base $k(e_1, \dots, e_d)$; the invariant-theoretic
    computation above is what effects the descent.
\end{rem*}

\subsubsection*{The degree count}

Throughout this paragraph, $\vec f$ is the reduced polynomial vector of
\Cref{lemma:realization_over_Gm_Wr:reorg}, so $m_j = \deg_{x^{-1}} f_j$ satisfies
$m_j < d\, p^{\,j-(r-1)}$, and $y_i := 1/x_i$, so that $f_j(x_i) \in y_i\, k[y_i]$ is
a polynomial in $y_i$ of degree $m_j$.

\begin{lemma}[symmetric regularity]\label{lemma:symmetric_regularity:reorg}
    Let $\Phi \in k[y_1, \dots, y_d]^{S_d}$ be symmetric with every monomial of total
    degree $< d$. Then $\Phi \in \kappa(\eta) = k(e_1/e_d, \dots, e_{d-1}/e_d)$; in
    particular $\Phi$ is regular at $\eta_\ndls$.
\end{lemma}

\begin{proof}
    $\Phi$ is a $k$-combination of monomial symmetric functions $m_\lambda(y)$ with
    $|\lambda| < d$. Each $m_\lambda(y)$ is an isobaric polynomial in the elementary
    symmetric functions $e_\mu(y) = \prod_i e_{\mu_i}(y)$ with $\mu \vdash |\lambda|$;
    every part $\mu_i \leq |\lambda| < d$, so only $e_1(y), \dots, e_{d-1}(y)$ occur
    --- never $e_d(y)$. Now $e_k(y) = e_{d-k}(x)/e_d(x)$, which for
    $1 \leq k \leq d-1$ lies in $\kappa(\eta)$. Hence
    $\Phi \in \kappa(\eta) \subset \widehat\Oo_\eta = \kappa(\eta)[[e_d]]$.
\end{proof}

\begin{theorem}[regularity of all Witt components]\label{theorem:witt_components_regularity:reorg}
    With $\vec f$ as in \Cref{lemma:realization_over_Gm_Wr:reorg}, every component
    $\alpha_n$ $(0 \leq n \leq r-1)$ of
    $\vec \alpha = \bigoplus_{i=1}^{d} \vec f(x_i)$ lies in $\kappa(\eta)$; in
    particular $\vec \alpha \in W_r\bigl(\widehat\Oo_\eta\bigr)$.
\end{theorem}

\begin{proof}
    Fix $n \leq r-1$ and a monomial $\prod_{i,j} \bigl(f_j(x_i)\bigr)^{e_{ij}}$ of
    $\alpha_n$, viewed as a universal isobaric polynomial in the $f_j(x_i)$ of
    weight $p^n$ (\Cref{lemma:witt_addition_isobaric:reorg}, with $f_j(x_i)$ of weight
    $p^j$); set $E_j := \sum_i e_{ij}$, so $\sum_j E_j p^{\,j} = p^{\,n}$. Expanding
    each $f_j(x_i)$ as a polynomial in $y_i$ of degree $m_j$, every resulting
    monomial in the $y_i$ has total degree
    \[
    \leq\; \sum_{i,j} e_{ij}\, m_j \;=\; \sum_j E_j\, m_j
    \;<\; \sum_j E_j\, d\, p^{\,j-(r-1)}
    \;=\; d\, p^{-(r-1)} \sum_j E_j p^{\,j}
    \;=\; d\, p^{\,n-(r-1)} \;\leq\; d,
    \]
    strictly: some $E_j > 0$ (as $\sum_j E_j p^j = p^n > 0$), and the pole bound
    $m_j < d\, p^{\,j-(r-1)}$ of
    \Cref{lemma:realization_over_Gm_Wr:reorg}\ref{item:realization_Wr_polebound:reorg} is strict
    even when $m_j = 0$. Thus every monomial of $\alpha_n$ --- a symmetric polynomial
    in $y_1, \dots, y_d$ by
    \Cref{prop:witt_cover_identification:reorg}\ref{item:witt_cover_symmetry:reorg} --- has total
    $y$-degree $< d$, and \Cref{lemma:symmetric_regularity:reorg} gives
    $\alpha_n \in \kappa(\eta)$.
\end{proof}

\begin{rem*}
    The bound is exact: the weight $p^{\,j-(r-1)}$ is paired against the isobaricity
    constraint $\sum_j E_j p^{\,j} = p^{\,n}$, so the estimate
    $d\, p^{\,n-(r-1)} \leq d$ is independent of how the Witt carries distribute the
    exponents $E_j$; the hypothesis ``break $< d$ at $P$'' is used at \emph{every}
    Witt level at once. The sanity check of
    \Cref{lemma:realization_over_Gm_Wr:reorg} records the extreme case: components with
    $p^{\,r-1-j} \geq d$ vanish outright in the reduced representative.
\end{rem*}

\subsubsection*{Proof of the theorem}

\begin{proof}[Proof of \Cref{theorem:ramification_after_blowup_is_tame_for_p_power:reorg}]
    By the standing reduction above we may assume $k$ algebraically closed and
    $\cP \to U$ totally ramified at $P$ with inertia all of $G$, so that
    $\cQ := \cP_x$ is a connected $\Z/p^r\Z$-torsor over $\Spec k((x))$ with
    ramification bounded by $d = \deg \ndls$.

    \Cref{lemma:realization_over_Gm_Wr:reorg} produces a torsor $\cP' \to \bG_m$, defined
    by a reduced polynomial vector $\vec f$ with $p^{r-1-j} m_j < d$, with
    $\cP'_x \cong \cQ$. By \Cref{cor:reduction_to_P1:reorg} (valid for arbitrary finite
    abelian $G$), the ramification of $\cP^{[d]}$ at $\eta_\ndls$ agrees with that of
    $\cP'^{[d]}$ at $\eta_{d \cdot 0}$; so we may assume $C = \bP^1_k$,
    $\mdls = d \cdot 0 + \infty$, $\ndls = d \cdot 0$, $\cP = \cP'$.

    By \Cref{prop:witt_cover_identification:reorg},
    \[
    \cP'^{[d]}\big|_{\Spec \widehat K_\eta}
    \;\cong\;
    \Spec\; \widehat K_\eta[\vec Z]\big/\bigl(\wp \vec Z \ominus \vec \alpha\bigr),
    \qquad
    \vec \alpha = \bigoplus_{i=1}^{d} \vec f(x_i),
    \]
    with $\widehat K_\eta = \kappa(\eta)((e_d))$ the complete local field at
    $\eta_\ndls$ (\Cref{eq:complete_field_at_generic_point:reorg}) and
    $\widehat\Oo_\eta = \kappa(\eta)[[e_d]]$ its valuation ring.

    By \Cref{theorem:witt_components_regularity:reorg},
    $\vec \alpha \in W_r(\kappa(\eta)) \subset W_r(\widehat\Oo_\eta)$. Hence, by
    \Cref{lemma:asw_etale_torsor:reorg} applied over $A = \widehat\Oo_\eta$, the
    Artin--Schreier--Witt cover $\wp \vec Z = \vec \alpha$ is finite \'etale over
    $\widehat\Oo_\eta$, and $\cP'^{[d]}\big|_{\Spec \widehat K_\eta}$ is its generic
    fibre: every field factor of
    $\widehat K_\eta[\vec Z]/(\wp \vec Z \ominus \vec \alpha)$ is the fraction field
    of a finite \'etale $\widehat\Oo_\eta$-algebra, i.e.\ an unramified extension of
    $\widehat K_\eta$ with respect to the discrete valuation ring
    $\widehat\Oo_\eta$. By \Cref{definition:tamely_ramified_in_codim_1:reorg},
    $\cP^{[d]}$ is unramified at $\eta_\ndls$.
\end{proof}

\begin{rem*}
    The argument proves all $r$ simultaneously; it reproves the case $r=1$
    (\Cref{cor:ramification_after_blowup_is_tame:reorg}(1)) as the special case of a
    vector of length one and never invokes it.
\end{rem*}


Finally, we assemble the general finite abelian case from the three cyclic
computations above.

\begin{proof}[Proof of \Cref{theorem:torsors_after_blowup_is_tame:reorg}]
    By the structure theorem, $G \cong \prod_i G_i$ with each $G_i$ cyclic of
    prime-power order $\ell_i^{a_i}$; write $\pi_i : G \to G_i$ for the
    projections. The $\pi_i$ jointly embed $G \hookrightarrow \prod_i G_i$, so
    the character $\rho : \Gal(\widehat K_{\eta_\ndls}^{sep}/\widehat K_{\eta_\ndls}) \to G$
    classifying $\cP^{[d]}$ at $\eta_\ndls$ is tamely ramified as soon as each
    component $\pi_i \circ \rho$ is: a homomorphism into $\prod_i G_i$ kills the
    wild inertia iff each of its projections does. Now $\pi_i \circ \rho$
    classifies the pushforward $\pi_{i,*}(\cP^{[d]})$, and the contracted
    product along $\pi_i$ commutes with the symmetric-power construction --- both
    are built from fibre products and finite quotients
    (\Cref{cor:sym-power-quotient}, \Cref{prop:torsor_module_equivalence}) ---
    so
    \[
        \pi_{i,*}\bigl(\cP^{[d]}\bigr) \cong \bigl(\pi_{i,*}\cP\bigr)^{[d]}.
    \]
    Here $\pi_{i,*}\cP$ is a $G_i$-torsor whose character is $\pi_i \circ \rho_{\cP}$,
    so its ramification is again bounded by $\mdls$ (composing with $\pi_i$ can
    only shrink the image of the ramification groups). Applying the cyclic cases:
    \begin{itemize}
        \item if $\ell_i \neq p$, then $\pi_{i,*}\cP$ is tamely ramified and
        $\bigl(\pi_{i,*}\cP\bigr)^{[d]}$ is tamely ramified at $\eta_\ndls$ by
        \Cref{cor:ramification_after_blowup_is_tame:reorg};
        \item if $\ell_i = p$, then $G_i \cong \Z/p^{a_i}\Z$ and
        $\bigl(\pi_{i,*}\cP\bigr)^{[d]}$ is \emph{unramified} at $\eta_\ndls$ by
        \Cref{theorem:ramification_after_blowup_is_tame_for_p_power:reorg}.
    \end{itemize}
    In either case $\pi_{i,*}(\cP^{[d]})$ is tamely ramified at $\eta_\ndls$;
    hence so is $\cP^{[d]}$.
\end{proof}




\section{Ramification on Products of Blowups}\label{subsection:ram_prod_analysis:reorg}
We conclude this sesction by proving sume auxilary results that would be useful in \Cref{section:gcft}.


%%% %%% %%% %%% %%% %%% %%% %%% %%% %%% %%% %%% %%% %%% %%% %%%
%       Behavior of Ramification under Product of Blowups
%%% %%% %%% %%% %%% %%% %%% %%% %%% %%% %%% %%% %%% %%% %%% %%%
%\subsubsection{Behavior of Ramification under Product of Blowups}

Let $X$ and $Y$ be smooth schemes over a field $k$, and let $x \in X$ and $y \in Y$ be closed points. 
We denote the blowups of these schemes at the given points by 
$\pi_X: \Bl_x(X) \to X$ and $\pi_Y: \Bl_y(Y) \to Y$. Furthermore, let $\pi_{X \times Y}: \Bl_{(x,y)}(X \times_k Y) \to X \times_k Y$ 
be the blowup of the product scheme at the point $(x,y)$.
We denote by $E_X, E_Y$, and $E_{X \times Y}$ the respective exceptional divisors, and let $\eta_X, \eta_Y$, and $\eta_{X \times Y}$ be their generic points.

In this section, we establish the following result concerning the stability of ramification bounds under the external product of torsors.

\begin{prop}\label{prop:ramification-external-product:reorg}
Let $G$ be a finite abelian group. Suppose $\cG_X$ and $\cG_Y$ are $G$-torsors defined on open subsets $U_X \subset \Bl_x(X)$ and $U_Y \subset \Bl_y(Y)$ that are disjoint 
from the exceptional divisors. 
If the ramification of $\cG_X$ at $\eta_X$ and $\cG_Y$ at $\eta_Y$ is bounded by $r$, then the external product torsors 
\[
\cG_{X \times Y} := \text{pr}_1^{-1} \cG_X \otimes \text{pr}_2^{-1} \cG_Y
\]
has ramification bounded by $r$ at the generic point $\eta_{X\times Y}$ of the exceptional divisor in the product blowup.
\end{prop}

The proposition is purely local in nature, it suffices to consider the case where $X$ and $Y$ are affine. 
More precisely, by the smoothness of $X$ and $Y$, we may restrict our attention to open neighborhoods of $x$ and $y$ that are etale over affine spaces. 
The rest of this section treats that case.

\subsection*{The Affine Case}
Let $X = \mathbb{A}_k^n$ be the affine $n$-space over a field $k$, and let $0 \in X$ be the origin. 
Let $\tilde{X} = \text{Bl}_0(X)$ be the blowup of $X$ at the origin. 
Recall that $\tilde{X} \subset X \times_k \mathbb{P}^{n-1}_k$ is defined by the equations $x_i u_j = x_j u_i$, where $[u_1 : \dots : u_n]$ are the homogeneous coordinates of 
$\mathbb{P}^{n-1}_k$.
The exceptional divisor $E \subset \tilde{X}$ is the fiber over the origin, $E = \{ (0, [u_1 : \dots : u_n]) \}$, which is of codimension 1 in $\tilde{X}$. 
Let $\eta \in E$ be the generic point of $E$, and let $R = \mathcal{O}_{\tilde{X}, \eta}$ be the associated local ring. 
This ring $R$ is a discrete valuation ring (DVR) with fraction field $K = K(X) = k(x_1, \dots, x_n)$.

On the affine chart $U_1$ where $u_1 \neq 0$, we have $x_i = \frac{u_i}{u_1} x_1$. 
The coordinate ring is:$$\mathcal{O}_{\tilde{X}}(U_1) = k\left[x_1, \frac{u_2}{u_1}, \dots, \frac{u_n}{u_1}\right]$$
In this chart, the generic point $\eta$ corresponds to the prime ideal $\mathfrak{p}_1 = (x_1)$. Thus, the local ring is $R = k[x_1, \frac{u_2}{u_1}, \dots, \frac{u_n}{u_1}]_{(x_1)}$. 
The residue field is $\kappa(\eta) = k(\frac{u_2}{u_1}, \dots, \frac{u_n}{u_1})$, and the completion of $R$ with respect to its maximal ideal is:
$$\hat{R} = \kappa(\eta)[[x_1]]$$
In this local ring, $x_1$ is a uniformizer. Note that any $x_i$ (for $i > 1$) can also serve as a uniformizer, 
as $x_i = (\frac{u_i}{u_1}) x_1$ and $\frac{u_i}{u_1}$ is a unit in $R$.

For any monomial $M = x_1^{a_1} \dots x_n^{a_n} \in K$, we can write:
$$M = x_1^{\sum a_i} \left( \frac{u_2}{u_1} \right)^{a_2} \dots \left( \frac{u_n}{u_1} \right)^{a_n}$$
Since the term in parentheses is a unit in $R$, the valuation $\nu_E$ associated with $E$ satisfies:
$$\nu_E(M) = \sum a_i = \deg M$$
Consequently, for any polynomial $f = f_d + f_{d+1} + \dots + f_l$, where $f_i$ is the homogeneous part of degree $i$, we have $\nu_E(f) = d$ (the order of vanishing at the origin).

\subsection*{The Product Case}
Now, let $X = \mathbb{A}^n$ and $Y = \mathbb{A}^m$ with origins $x=0$ and $y=0$. 
As before, $\text{Bl}_0(X) \subset X \times \mathbb{P}^{n-1}$ and $\text{Bl}_0(Y) \subset Y \times \mathbb{P}^{m-1}$ have exceptional divisors $E_X$ and $E_Y$ respectively.
Consider the product $X \times_k Y \cong \mathbb{A}_k^{n+m}$. 
The blowup of the product at the origin $(0,0)$, denoted $\text{Bl}_{(0,0)}(X \times_k Y)$, is a subscheme of $(X \times Y) \times \mathbb{P}^{n+m-1}$ defined by:

$$\begin{cases} 
x_i w_j = x_j w_i & 1 \le i, j \le n \\
y_k w_{n+l} = y_l w_{n+k} & 1 \le k, l \le m \\
x_i w_{n+j} = y_j w_i & 1 \le i \le n, \,\, 1 \le j \le m 
\end{cases}$$
where $[w_1 : \dots : w_{n+m}]$ are the homogeneous coordinates of $\mathbb{P}^{n+m-1}$. The exceptional divisor $E_{X \times Y}$ is isomorphic to $\mathbb{P}^{n+m-1}$.

\subsection*{Comparison of Blowups}

Both $\text{Bl}_0(X) \times_k \text{Bl}_0(Y)$ and $\text{Bl}_{(0,0)}(X \times_k Y)$ are birational to $X \times_k Y$. 
They share a common dense open set $\tilde{U}$ defined by the condition that neither the $X$-coordinates nor the $Y$-coordinates vanish simultaneously in the projective space:
$$\tilde{U} = \{ ((x,y), [w_1: \dots : w_{n+m}]) \mid (w_1, \dots, w_n) \neq 0 \text{ and } (w_{n+1}, \dots, w_{n+m}) \neq 0 \}$$This yields a diagram of open immersions:
$$ 
   \begin{tikzcd}
        & \tilde{U} \arrow[dl, "f_1"'] \arrow[dr, "f_2", hook] & \\
        \text{Bl}_0(X) \times_k \text{Bl}_0(Y) & & \text{Bl}_{(0,0)}(X \times_k Y)
    \end{tikzcd}
$$
    
    where $f_1$ maps the coordinates to the respective projectivizations $[w_1: \dots : w_n]$ and $[w_{n+1}: \dots : w_{n+m}]$.
    Also note that the generic point $\eta_{X \times Y}$  of $E_{X \times_k Y}$ is inside $\tilde{U}$

\subsection*{Extensions of DVRs}
Let $S$ be the local ring of the generic point $\eta_{X\times Y}$ of $E_{X \times Y}$ in $\tilde{U}$. 
On the chart where $w_1 \neq 0$ and $w_{n+1} \neq 0$, we have 
$x_1 = (\frac{w_1}{w_{n+1}}) y_1$. 
Since $\frac{w_1}{w_{n+1}}$ is a unit in this chart, $x_1$ and $y_1$ are equivalent as uniformizers.
We have:
$$
S=k\left[x_1, \frac{w_2}{w_1} \cdots, \frac{w_n}{w_1}, \frac{w_{n+1}}{w_1} \cdots \frac{w_{n+m}}{w_1}\right]_{(x_1)} = 
k\left[\frac{w_1}{w_{n+1}}, \frac{w_2}{w_{n+1}} \cdots \frac{w_n}{w_{n+1}}, y_1 \cdots \frac{w_{n+m}}{w_{n+1}}\right]_{(y_1)}
$$

$$k(\eta_{X \times Y}) = k\left(\frac{w_2}{w_1} \cdots, \frac{w_n}{w_1}, \frac{w_{n+1}}{w_1} \cdots \frac{w_{n+m}}{w_1}\right)$$

Let $R_X$ be the local ring of the exceptional divisor $E_X$ in $\text{Bl}_0(X)$. 
The pullback of $E_X \times_k \text{Bl}_0(Y)$ along $f_1$ induces an extension of DVRs $R_X \hookrightarrow S$. Which is:
\begin{enumerate}
    \item Weakly Unramified: $x_1$ is a uniformizer in both $R_X$ and $S$, so the ramification index is $e=1$.
    \item Residually Transcendental: The residue field extension $\kappa(\eta_X) \subset \kappa(\eta)$ is:
    $$k\left(\frac{w_2}{w_1}, \dots, \frac{w_n}{w_1}\right) \subset k\left(\frac{w_2}{w_1}, \dots, \frac{w_n}{w_1}, \frac{w_{n+1}}{w_1}, \dots, \frac{w_{n+m}}{w_1}\right)$$
    Hence separable.
\end{enumerate}
Since this extension is generated by transcendental elements, it is separable and formally smooth at the maximal ideal (\stackstag{09E7}).

\subsection*{Ramification of $G$-Torsors}
Let $G$ be a finite abelian group. 
Let $\mathcal{Q}$ be a $G$-torsor on an open $U \subset \text{Bl}_0(X)$ disjoint from $E_X$. 
%Suppose $\mathcal{Q}$ has ramification bounded by $r$ at the generic point $\eta_X$ of $E_X$.
Let $V \subset \text{Bl}_0(Y)$ be an open subscheme, and let $\pi_X: \text{Bl}_0(X) \times_k V \to \text{Bl}_0(X)$ 
be the projection onto the first factor. 
By restricting this projection to $U \times_k V$, we obtain the pullback $G$-torsor:
$$\pi_X^{-1}(\mathcal{Q}) \cong \mathcal{Q} \times_k V$$
which is defined on the open subset $U \times_k V \subset \text{Bl}_0(X) \times_k \text{Bl}_0(Y)$.
The extension of local rings $R_X \to S$ is weakly unramified (the ramification index $e=1$) and residually transcendental with seprable residue field exntesion. 
Under these conditions the ramification filtration is preserved. 
Therefore, the pullback $\pi_X^{-1}(\mathcal{Q})$ has ramification bounded by $r$ at the generic point $\eta_{X \times Y}$ of the exceptional divisor $E_{X \times Y}$ if and only if 
the original torsor $\mathcal{Q}$ has ramification bounded by $r$ at the generic point $\eta_X$ of $E_X$

And we finish by \Cref{lemma:bounding_ramification_of_contracted_product:reorg}.

